\chapter{量子计算基础}
\label{ch:lec01}

\begin{objectives}
\begin{itemize}[leftmargin=1.4em,nosep]
  \item 用量子比特与寄存器的语言陈述量子力学的公设。
  \item 通过张量积与狄拉克记号操作多量子比特态。
  \item 从通用门组读懂并构建简单的量子线路。
  \item 用 CNOT 门与 Toffoli 门构建可逆的二进制加法器。
  \item 通过门合成制备态，并描述投影测量与广义（POVM）测量。
  \item 逐个振幅地追踪量子隐形传态与超密编码协议。
  \item 解释测量、不可克隆定理、不可超光速通信以及密度算符。
  \item 通过 CHSH 不等式及其量子值 $2\sqrt2$ 量化非定域性。
  \item 定义量子算法的基本复杂度度量。
\end{itemize}
\end{objectives}

\coverage{\S1.1、\S1.2（不含 \S1.2.4）、\S1.3、\S1.4、\S1.5.1 以及 \S1.6。本学期不
涵盖：\S1.2.4（非定域性与 CHSH 不等式）、\S1.5.2（超密编码）、\S1.7（密度算符与混合
态）以及 \S1.8（量子算法的复杂度）。}

本讲构建起整个量子计算所依托的数学舞台。我们所需的唯一公设看似简单得出人意料：
\emph{一个量子系统的态是复内积空间中的单位向量。}由这一个想法出发，我们便可
重建狄拉克的左矢--右矢记号、交叠与正交的几何、单个量子比特的描述与布洛赫球，并
且——借助张量积——重建多个量子比特那指数级庞大的态空间，纠缠正是在此首次出现。
随后我们转向操作层面：作为酉算符的门、标准的单量子比特门与多量子比特门、量子线路、
通用门组，以及不可克隆定理。我们以两个贯穿全课程的想法收尾：密度算符形式，它让我们
得以描述子系统与混合；以及我们用以评判量子算法的复杂度度量（查询复杂度、门数、线路
深度）。后续各讲的一切都建立在这套词汇之上。

\section{量子比特、叠加与测量}
\label{sec:lec01-1}

我们从线性代数的框架开始。

\begin{definition}[复内积空间]\label{def:lec01-inner}
设 $\Hilb$ 为一复向量空间。$\Hilb$ 上的一个\emph{内积}是一个映射
$\langle\,\cdot\,,\,\cdot\,\rangle:\Hilb\times\Hilb\to\C$，对所有
$u,v,w\in\Hilb$ 与 $a,b\in\C$ 满足：
\begin{enumerate}[label=(\roman*)]
  \item \emph{共轭对称性：} $\langle u,v\rangle=\overline{\langle v,u\rangle}$。
  \item \emph{第二变元线性：}
        $\langle u,\,a v+b w\rangle=a\langle u,v\rangle+b\langle u,w\rangle$。
        （与 (i) 结合，这使得内积在第一变元上\emph{共轭线性}。）
  \item \emph{正定性：} $\langle u,u\rangle\ge 0$，等号成立当且仅当 $u=0$。
\end{enumerate}
\end{definition}

我们采用物理学家的约定：内积在\emph{第二}个槽位上线性、在第一个槽位上共轭线性；
正是这一约定使狄拉克记号得以干净利落地呈现。具体地，在 $\C^{N}$ 上，取
$x=(x_0,\dots,x_{N-1})^{T}$ 与 $y=(y_0,\dots,y_{N-1})^{T}$，
\begin{equation}
  \langle x,y\rangle \;=\; x\dg y \;=\; \sum_{i=0}^{N-1}\overline{x_i}\,y_i .
\end{equation}

\begin{definition}[范数与单位向量]\label{def:lec01-norm}
由内积诱导的\emph{范数}为 $\norm{u}=\sqrt{\langle u,u\rangle}$。
若 $\norm{u}=1$，则向量 $u$ 称为\emph{单位向量}（或\emph{已归一化的}）。两个
向量 $u,v$ 若 $\langle u,v\rangle=0$，则称它们\emph{正交}。
\end{definition}

\begin{definition}[有限维希尔伯特空间]\label{def:lec01-hilbert}
一个\emph{有限维希尔伯特空间 (Hilbert space)} 是一个有限维复
内积空间 $\Hilb$（在有限维情形完备性自动成立，故无需额外条件）。若 $\dim\Hilb=N$，则选定一组标准正交基
便给出一个保持内积的同构 $\Hilb\cong\C^{N}$。对于
由 $n$ 个量子比特构成的系统，我们将得到 $N=2^{n}$。
\end{definition}

现在可以精确地陈述态空间公设了。

\begin{remark}[态空间公设]
一个孤立量子系统的态由希尔伯特空间 $\Hilb$ 中的一个\emph{单位}向量描述。该
空间对线性组合封闭：若
$u,v\in\Hilb$ 且 $a,b\in\C$，则 $a u+b v\in\Hilb$。这种封闭性即为
\emph{叠加原理}——任何合法态的线性组合仍是一个合法的（未归一化的）态。
\end{remark}

\subsection{内积主宰一切几何}

\begin{proposition}[公理的初等推论]\label{prop:lec01-elem}
对所有 $u,v\in\Hilb$ 与 $a\in\C$：
\begin{enumerate}[label=(\alph*)]
  \item $\langle u,0\rangle=\langle 0,u\rangle=0$；
  \item $\norm{a u}=\abs{a}\,\norm{u}$；
  \item \emph{（勾股定理）}若 $\langle u,v\rangle=0$，则
        $\norm{u+v}^{2}=\norm{u}^{2}+\norm{v}^{2}$。
\end{enumerate}
\end{proposition}

\begin{proof}
(a) 由第二变元线性，$\langle u,0\rangle=\langle u,0\cdot 0\rangle=0\,\langle u,0\rangle=0$；
另一个等式由共轭对称性得到。
(b) $\norm{a u}^{2}=\langle a u,a u\rangle=\overline{a}\,a\,\langle u,u\rangle=\abs{a}^{2}\norm{u}^{2}$，
再取平方根。
(c) 展开 $\norm{u+v}^{2}=\langle u+v,u+v\rangle=\norm{u}^{2}+\langle u,v\rangle+\langle v,u\rangle+\norm{v}^{2}$；
两个交叉项消失，因为 $\langle u,v\rangle=0$ 迫使
$\langle v,u\rangle=\overline{\langle u,v\rangle}=0$。
\end{proof}

整个学科中最有用的单个不等式便是下面这个。

\begin{theorem}[柯西--施瓦茨不等式]\label{thm:lec01-cs}
对所有 $u,v\in\Hilb$，
\begin{equation}
  \abs{\langle u,v\rangle}\;\le\;\norm{u}\,\norm{v},
\end{equation}
等号成立当且仅当 $u$ 与 $v$ 线性相关。
\end{theorem}

\begin{proof}
若 $v=0$，则两边皆为 $0$ 且 $u,v$ 平凡地相关，故设 $v\neq0$。
对任意 $\lambda\in\C$，正定性给出
\[
  0\;\le\;\norm{u-\lambda v}^{2}
   \;=\;\norm{u}^{2}-\overline{\lambda}\langle v,u\rangle-\lambda\langle u,v\rangle+\abs{\lambda}^{2}\norm{v}^{2}.
\]
取 $\lambda=\dfrac{\langle v,u\rangle}{\norm{v}^{2}}$（使
右端取极小的值）。于是 $\lambda\langle u,v\rangle=\dfrac{\langle v,u\rangle\langle u,v\rangle}{\norm{v}^{2}}=\dfrac{\abs{\langle u,v\rangle}^{2}}{\norm{v}^{2}}$，同样地 $\overline{\lambda}\langle v,u\rangle=\dfrac{\abs{\langle u,v\rangle}^{2}}{\norm{v}^{2}}$，故
\[
  0\;\le\;\norm{u}^{2}-\frac{\abs{\langle u,v\rangle}^{2}}{\norm{v}^{2}}.
\]
整理即得 $\abs{\langle u,v\rangle}^{2}\le\norm{u}^{2}\norm{v}^{2}$，亦即
所证。等号成立当且仅当 $\norm{u-\lambda v}=0$，即 $u=\lambda v$，这
恰为线性相关。
\end{proof}

\begin{corollary}[三角不等式]\label{cor:lec01-tri}
对所有 $u,v\in\Hilb$，$\norm{u+v}\le\norm{u}+\norm{v}$。
\end{corollary}

\begin{proof}
$\norm{u+v}^{2}=\norm{u}^{2}+2\,\mathrm{Re}\,\langle u,v\rangle+\norm{v}^{2}
 \le \norm{u}^{2}+2\abs{\langle u,v\rangle}+\norm{v}^{2}
 \le \norm{u}^{2}+2\norm{u}\norm{v}+\norm{v}^{2}=(\norm{u}+\norm{v})^{2}$，
这里用到 $\mathrm{Re}\,z\le\abs{z}$ 以及 \cref{thm:lec01-cs}。取平方根即可。
\end{proof}

对单位向量而言，柯西--施瓦茨不等式即为 $\abs{\langle u,v\rangle}\le1$，
故可写作
\begin{equation}\label{eq:lec01-geomdecomp}
  \langle u,v\rangle \;=\; \underbrace{\norm{u}\,\norm{v}}_{\text{length}}\;
  \underbrace{\cos\theta}_{\text{alignment}}\;\underbrace{\ee^{\ii\varphi}}_{\text{relative phase}},
  \qquad \cos\theta=\frac{\abs{\langle u,v\rangle}}{\norm{u}\,\norm{v}}\in[0,1].
\end{equation}
模量度两个方向重合的程度（$\cos\theta=1$ 表示
平行，$\cos\theta=0$ 表示正交），而辐角 $\varphi$ 则是一个
\emph{相对相位}，它对实向量没有对应物。正交性，
即 $\langle u,v\rangle=0$，恰好表述了两个态完全
可区分。

\subsection{狄拉克记号}

\begin{definition}[右矢与左矢]\label{def:lec01-braket}
$\Hilb$ 的一个向量写作\emph{右矢} $\ket{x}$。每个右矢都
对应一个\emph{左矢} $\bra{x}$，它是对偶空间 $\Hilb^{*}$ 的元素，
以 $\bra{x}:\ket{y}\mapsto\langle x,y\rangle$ 作用。在坐标下，若 $\ket{x}$
是以 $x_i$ 为分量的列向量，则 $\bra{x}$ 是行向量
$x\dg=(\overline{x_0},\dots,\overline{x_{N-1}})$，即其共轭转置。
\end{definition}

\begin{definition}[内积与外积]\label{def:lec01-inout}
$\ket{x}$ 与 $\ket{y}$ 的\emph{内积}是标量
$\braket{x|y}=\langle x,y\rangle\in\C$。\emph{外积} $\ket{x}\!\bra{y}$
是如下作用的线性算符
\[
  \bigl(\ket{x}\!\bra{y}\bigr)\ket{z}=\ket{x}\,\braket{y|z}=\braket{y|z}\,\ket{x}.
\]
在坐标下它是以 $\bigl(\ket{x}\!\bra{y}\bigr)_{ij}=x_i\,\overline{y_j}$ 为分量的矩阵。
\end{definition}

固定一组标准正交基 $\{\ket{i}\}_{i=0}^{N-1}$，即 $\braket{i|j}=\delta_{ij}$。
任何右矢都展开为 $\ket{x}=\sum_i x_i\ket{i}$，而对其作用 $\bra{j}$ 则挑出
单个系数：
\begin{equation}\label{eq:lec01-coeff}
  \braket{j|x}=\sum_i x_i\braket{j|i}=\sum_i x_i\,\delta_{ji}=x_j .
\end{equation}
于是第 $j$ 个振幅就是交叠 $\braket{j|x}$。

\begin{theorem}[完备性 / 单位分解]\label{thm:lec01-completeness}
若 $\{\ket{i}\}_{i=0}^{N-1}$ 是 $\Hilb$ 的一组标准正交基，则
\begin{equation}
  \sum_{i=0}^{N-1}\proj{i}\;=\;\id .
\end{equation}
因此每个算符 $A$ 都满足 $A=\sum_{i,j}A_{ij}\,\ket{i}\!\bra{j}$，其中
$A_{ij}=\braket{i|A|j}$。
\end{theorem}

\begin{proof}
令 $P=\sum_i\proj{i}$。对任意 $\ket{x}=\sum_j x_j\ket{j}$，
\[
  P\ket{x}=\sum_i\ket{i}\braket{i|x}=\sum_i x_i\ket{i}=\ket{x},
\]
这里用到 \eqref{eq:lec01-coeff}。由于对每个 $\ket{x}$ 都有 $P\ket{x}=\ket{x}$，故
$P=\id$。至于第二个结论，在 $A$ 两侧插入单位算符：
$A=\id\,A\,\id=\sum_{i,j}\ket{i}\braket{i|A|j}\bra{j}=\sum_{i,j}A_{ij}\,\ket{i}\!\bra{j}$。
\end{proof}

\begin{remark}[投影算符]
对单位向量 $\ket{u}$，算符 $P_u=\proj{u}$ 是到由 $\ket{u}$ 张成的直线上的正交
\emph{投影算符}：它是厄米的
（$P_u\dg=P_u$）且幂等的，因为 $P_u^{2}=\ket{u}\braket{u|u}\bra{u}=\proj{u}=P_u$，
这是由于 $\braket{u|u}=1$。\cref{thm:lec01-completeness} 的完备性关系
把 $\id$ 分解为若干正交的秩一
投影算符之和，每个基向量对应一个。
\end{remark}

\subsection{单个量子比特}

\begin{definition}[量子比特与计算基]\label{def:lec01-qubit}
一个\emph{量子比特}是态空间为 $\Hilb=\C^{2}$ 的量子系统。那组
特别指定的标准正交基
\[
  \ket{0}=\begin{pmatrix}1\\0\end{pmatrix},\qquad
  \ket{1}=\begin{pmatrix}0\\1\end{pmatrix}
\]
称为\emph{计算基}。一般的纯态为
\begin{equation}\label{eq:lec01-qubit}
  \ket{\psi}=a\ket{0}+b\ket{1},\qquad a,b\in\C,\qquad \abs{a}^{2}+\abs{b}^{2}=1 .
\end{equation}
系数 $a,b$ 称为\emph{振幅}；$\abs{a}^{2}$ 与 $\abs{b}^{2}$
是在计算基测量中得到 $0$ 或 $1$ 的概率。
\end{definition}

刚才陈述的法则——结果 $i$ 以概率 $\abs{\braket{i|\psi}}^{2}$ 出现，
且测量后的态坍缩到 $\ket{i}$——即为\emph{玻恩规则}，
量子力学的测量公设。更一般地，一个带有正交投影算符 $\{P_m\}$（$\sum_m P_m=\id$）的投影
测量以概率 $p_m=\bra{\psi}P_m\ket{\psi}$ 返回结果
$m$，并使态变为 $P_m\ket{\psi}/\sqrt{p_m}$；由完备性，这些概率之和为一。

\begin{remark}[整体相位不可观测]\label{rem:lec01-global}
测量概率仅通过形如 $\abs{\braket{i|\psi}}^{2}$ 的模依赖于振幅。
对一个实的\emph{整体相位} $\gamma$，用 $\ee^{\ii\gamma}\ket{\psi}$ 替换 $\ket{\psi}$ 会
使每一个这样的概率保持不变，更
一般地，当同一相位出现在两侧时，它会使 $\braket{\phi|A|\psi}$ 型的量保持不变。
因此 $\ee^{\ii\gamma}\ket{\psi}$ 与 $\ket{\psi}$ 描述
同一个物理态：真正的态是\emph{射影空间}
$\Prob(\Hilb)$ 中的点，即射线而非向量。叠加各项之间的\emph{相对}相位则
恰恰相反，是物理的，并驱动干涉。
\end{remark}

\begin{proposition}[布洛赫参数化]\label{prop:lec01-bloch}
每个单量子比特纯态，在相差一个不可观测的整体相位下，都可写为
\begin{equation}\label{eq:lec01-blochform}
  \ket{\psi}=\cos\tfrac{\theta}{2}\,\ket{0}+\ee^{\ii\varphi}\sin\tfrac{\theta}{2}\,\ket{1},
  \qquad \theta\in[0,\pi],\quad \varphi\in[0,2\pi).
\end{equation}
数对 $(\theta,\varphi)$ 是单位
球面 $S^{2}\subset\R^{3}$ 上一点的球面角对，该球面即\emph{布洛赫球}。
\end{proposition}

\begin{proof}
从带有 $\abs{a}^{2}+\abs{b}^{2}=1$ 的 \eqref{eq:lec01-qubit} 出发。把
振幅写成极坐标形式 $a=r_a \ee^{\ii\alpha}$ 与 $b=r_b \ee^{\ii\beta}$，其中
$r_a,r_b\ge0$。由于 $r_a^{2}+r_b^{2}=1$，我们可对唯一的 $\theta\in[0,\pi]$ 置 $r_a=\cos\tfrac{\theta}{2}$ 与
$r_b=\sin\tfrac{\theta}{2}$（此时 $\theta/2$ 的余弦与正弦
皆非负）。于是
\[
  \ket{\psi}=\ee^{\ii\alpha}\cos\tfrac{\theta}{2}\ket{0}+\ee^{\ii\beta}\sin\tfrac{\theta}{2}\ket{1}
           =\ee^{\ii\alpha}\Bigl(\cos\tfrac{\theta}{2}\ket{0}+\ee^{\ii(\beta-\alpha)}\sin\tfrac{\theta}{2}\ket{1}\Bigr).
\]
整体因子 $\ee^{\ii\alpha}$ 是一个整体相位，由
\cref{rem:lec01-global} 可将其舍弃。令 $\varphi=\beta-\alpha\ (\mathrm{mod}\ 2\pi)$
便得 \eqref{eq:lec01-blochform}。当 $\theta\in\{0,\pi\}$ 时态为 $\ket{0}$
或 $\ket{1}$，而 $\varphi$ 无关紧要（球面的两极），这与
态同 $S^{2}$ 上的点的对应是一致的。
\end{proof}

\begin{remark}[读懂布洛赫球]
在这幅图景中，$\ket{0}$ 与 $\ket{1}$ 位于北极与南极
（$\theta=0,\pi$），而\emph{正交的}量子比特态是\emph{对径}点，
而非垂直的点——一个态与其正交补在球面上相隔
角度 $\theta=\pi$，尽管它们在 $\C^{2}$ 中相隔 $90^{\circ}$。赤道上的态 $\tfrac{1}{\sqrt2}(\ket{0}+\ee^{\ii\varphi}\ket{1})$
具有 $\theta=\pi/2$，彼此仅在相对相位 $\varphi$ 上不同。
\end{remark}

\begin{example}[振幅与归一化]\label{ex:lec01-pm}
态 $\ket{+}=\tfrac{1}{\sqrt2}(\ket{0}+\ket{1})$ 具有
$\abs{a}^{2}=\abs{b}^{2}=\tfrac12$，故计算基测量以相等的概率 $\tfrac12$ 给出 $0$
或 $1$；其布洛赫角为 $\theta=\pi/2,\varphi=0$。
态 $\ket{-}=\tfrac{1}{\sqrt2}(\ket{0}-\ket{1})$ 在计算基中具有\emph{相同}的测量
统计，但相差一个相对相位
$\varphi=\pi$；事实上 $\braket{+|-}=\tfrac12(1-1)=0$，故 $\ket{+}$ 与 $\ket{-}$
正交，并可由适当的测量完全区分。
\end{example}
\section{多量子比特系统与张量积}
\label{sec:lec01-2}

\begin{definition}[张量积态空间]\label{def:lec01-tensor}
若系统 $A$ 与 $B$ 的态空间分别为 $\Hilb_A$ 与 $\Hilb_B$，则复合系统的态空间为
\emph{张量积 (tensor product)} $\Hilb_A\otimes\Hilb_B$。若
$\{\ket{i}_A\}$ 与 $\{\ket{j}_B\}$ 是标准正交基 (orthonormal basis)，则
$\{\ket{i}_A\otimes\ket{j}_B\}$ 是该积空间的标准正交基，其内积 (inner product) 为
\[
  \bigl(\bra{i}\otimes\bra{j}\bigr)\bigl(\ket{k}\otimes\ket{l}\bigr)
    =\braket{i|k}\,\braket{j|l}.
\]
因此 $\dim(\Hilb_A\otimes\Hilb_B)=(\dim\Hilb_A)(\dim\Hilb_B)$。对 $n$ 个量子比特 (qubit)，
\[
  \Hilb=(\C^{2})^{\otimes n}\cong\C^{2^{n}},\qquad \dim\Hilb=2^{n}.
\]
\end{definition}

我们把 $\ket{i}\otimes\ket{j}$ 简记为 $\ket{i}\ket{j}$ 或 $\ket{ij}$。$n$ 个量子比特的
计算基由 $2^{n}$ 个比特串索引，即 $\{\ket{x}:x\in\{0,1\}^{n}\}$。维数的指数增长——$50$ 个
量子比特就已给出维数为 $2^{50}>10^{15}$ 的空间——是量子计算潜在威力的一个根源：对一般态的
经典描述需要指数多个振幅 (amplitude)。

\begin{remark}[计算基态的 $\ket{k}$ 记号]\label{rem:lec01-ket-k}
将整数 $k$ 写入右矢 (ket) 表示“独热 (one-hot)”向量，它在索引 $k$ 处有唯一的 $1$，其余处为 $0$。
于是对两个量子比特而言
\[
  \ket{00}=\begin{psmallmatrix}1\\0\\0\\0\end{psmallmatrix}=\ket{0}_4,\quad
  \ket{01}=\ket{1}_4,\quad \ket{10}=\ket{2}_4,\quad \ket{11}=\ket{3}_4,
\]
即四个独热的 $4$-维向量。注意 $\ket{00}=\ket{0}_4$ \emph{不是}单量子比特的
$\ket{0}_2$；当维数在上下文中不明确时，我们用 $\ket{k}_d$ 标注。一般地，
$\ket{x_1\cdots x_n}=\ket{k}$，其中 $k\in\{0,\dots,2^{n}-1\}$ 是以 $x_1\cdots x_n$ 为
二进制表示的整数；这些独热态恰好就是计算基，由 $\braket{k|m}=\delta_{km}$ 可知其标准正交。
\end{remark}

\subsection{乘积态与纠缠态}
\label{sec:lec01-buildprod}

\begin{definition}[乘积态与纠缠态]\label{def:lec01-product}
若态 $\ket{\psi}\in\Hilb_A\otimes\Hilb_B$ 能写成
$\ket{\psi}=\ket{\alpha}_A\otimes\ket{\beta}_B$（对某个
$\ket{\alpha}_A\in\Hilb_A$、$\ket{\beta}_B\in\Hilb_B$），则称 $\ket{\psi}$ 为
\emph{乘积态 (product state)}（或\emph{可分的 (separable)}）。否则 $\ket{\psi}$ 是
\emph{纠缠的 (entangled)}。
\end{definition}

为了解乘积态是如何组装起来的，不妨亲手计算一个张量积，因为其结果不像把两个向量叠放那样显然。

\begin{example}[两个量子比特的张量积，逐分量计算]\label{exa:lec01-tp-explicit}
设 $\ket{\psi}=\psi_0\ket0+\psi_1\ket1=\begin{psmallmatrix}\psi_0\\\psi_1\end{psmallmatrix}$，
$\ket{\varphi}=\begin{psmallmatrix}\varphi_0\\\varphi_1\end{psmallmatrix}$。把 $\ket{\psi}$ 的
\emph{每个}分量都与 $\ket{\varphi}$ 的\emph{全部}配对：
\begin{equation}\label{eq:lec01-tp}
  \ket{\psi}\otimes\ket{\varphi}
   =\begin{pmatrix}\psi_0\begin{psmallmatrix}\varphi_0\\\varphi_1\end{psmallmatrix}\\[6pt]
                   \psi_1\begin{psmallmatrix}\varphi_0\\\varphi_1\end{psmallmatrix}\end{pmatrix}
   =\begin{pmatrix}\psi_0\varphi_0\\ \psi_0\varphi_1\\ \psi_1\varphi_0\\ \psi_1\varphi_1\end{pmatrix}.
\end{equation}
这四个振幅就是从每个量子比特各取一个振幅的四个\emph{乘积}，依次排列；结果是一个描述两个量子
比特整体的 $4$-分量向量。它的范数得以保持，
$\sum_{i,j}\abs{\psi_i\varphi_j}^{2}=(\abs{\psi_0}^2+\abs{\psi_1}^2)(\abs{\varphi_0}^2+\abs{\varphi_1}^2)=1$，
故它仍是一个合法的态。\textbf{警告：}$\ket{\psi}\otimes\ket{\varphi}$ \emph{不是}把 $\ket{\psi}$
的两个分量叠放在 $\ket{\varphi}$ 的分量之上——它是 \eqref{eq:lec01-tp} 中逐分量的乘积。
\end{example}

\begin{example}[用门构建的乘积态]\label{exa:lec01-prod-gates}
让两个量子比特都从 $\ket0$ 出发，并作用单量子比特门 (gate)（上方的量子比特写在左边；线路中的
“自上而下”对应右矢中的“从左到右”）：
\[
  \begin{aligned}
   \ket0\otimes\ket0 &= \ket{00}=\ket{0}_4, \\
   (X\ket0)\otimes\ket0 &= \ket1\otimes\ket0=\ket{10}=\ket{2}_4, \\
   \ket0\otimes(H\ket0) &= \ket0\otimes\ket{+}=\tfrac1{\sqrt2}(\ket{00}+\ket{01}), \\
   (H\ket0)\otimes(H\ket0) &= \ket{+}\otimes\ket{+}
       =\tfrac12(\ket{00}+\ket{01}+\ket{10}+\ket{11})=\tfrac12\sum_{k=0}^{3}\ket{k}.
  \end{aligned}
\]
最后一行——四个振幅都等于 $\tfrac12$——是均匀叠加 (superposition)，其 $n$ 量子比特版本
$H^{\otimes n}\ket{0\cdots0}$ 是许多量子算法的标准开场动作（\cref{prop:lec01-prepbasic}）。
这里的每个态都是乘积态；有意思的问题是：一个两量子比特态何时\emph{不能}以这种方式构建——这正是
接下来几个小节的主题。
\end{example}

\begin{example}[一个一般的两量子比特态]
两个量子比特处于 $\C^{4}$ 中，其基为 $\ket{00},\ket{01},\ket{10},\ket{11}$。一个一般的态为
\[
  \ket{\psi}=\sum_{i,j\in\{0,1\}}x_{ij}\ket{ij},\qquad
  \sum_{i,j}\abs{x_{ij}}^{2}=1 .
\]
例如 $\tfrac12(\ket{00}+\ket{01}+\ket{10}+\ket{11})$ 是归一化的
（$4\cdot\tfrac14=1$），并且，读者可以验证，它分解为
$\tfrac{1}{\sqrt2}(\ket{0}+\ket{1})\otimes\tfrac{1}{\sqrt2}(\ket{0}+\ket{1})$，故它是乘积态。
\end{example}

\subsection{检测纠缠：可分性判据}
\label{sec:lec01-septest-sub}

通过把振幅看作一个矩阵，我们可以具体地判定可分性，而无需去猜测某个分解。

\begin{proposition}[两量子比特的可分性判据]\label{prop:lec01-septest}
将两量子比特态写成 $\ket{\psi}=\sum_{i,j\in\{0,1\}}x_{ij}\ket{ij}$，并把其振幅收入
$2\times2$ 的\emph{系数矩阵} $X=\begin{psmallmatrix}x_{00}&x_{01}\\ x_{10}&x_{11}\end{psmallmatrix}$。
则 $\ket{\psi}$ 是乘积态当且仅当
\[
  \det X \;=\; x_{00}x_{11}-x_{01}x_{10}\;=\;0
  \qquad(\text{等价地，} \rank X=1).
\]
\end{proposition}

\begin{proof}
若 $\ket{\psi}=\ket{\alpha}\otimes\ket{\beta}$，其中 $\ket{\alpha}=a_0\ket0+a_1\ket1$、
$\ket{\beta}=b_0\ket0+b_1\ket1$，则 $x_{ij}=a_i b_j$，即
$X=\mathbf{a}\,\mathbf{b}^{\mathsf T}$ 是一个外积；因此 $\rank X\le1$ 且
$\det X=a_0b_0\,a_1b_1-a_0b_1\,a_1b_0=0$。反之，一个 $\det X=0$ 的非零 $2\times2$ 矩阵的秩为
$1$，故它可对某些向量 $\mathbf a,\mathbf b$ 分解为 $X=\mathbf{a}\,\mathbf{b}^{\mathsf T}$；
读出 $x_{ij}=a_ib_j$ 便得 $\ket{\psi}=\ket{\alpha}\otimes\ket{\beta}$（整体归一化可在两个
因子之间分配）。更一般地，对一个二分态而言，$X$ 的非零奇异值个数即\emph{施密特秩 (Schmidt rank)}：
秩为 $1$ 意味着乘积态，秩 $>1$ 意味着纠缠。
\end{proof}

\begin{example}[判据的实际运用]\label{exa:lec01-septest}
\begin{itemize}[nosep,leftmargin=1.4em]
  \item $\ket{+}\ket{0}=\tfrac1{\sqrt2}(\ket{00}+\ket{10})$ 有
        $X=\tfrac1{\sqrt2}\begin{psmallmatrix}1&0\\ 1&0\end{psmallmatrix}$，故 $\det X=0$：
        是乘积态（正如其名所示）。
  \item $\tfrac12(\ket{00}+\ket{01}+\ket{10}+\ket{11})$ 的每个 $x_{ij}=\tfrac12$，故
        $\det X=\tfrac14-\tfrac14=0$：即上面的乘积态 $\ket{+}\ket{+}$。
  \item 贝尔态 (Bell state) $\ket{\Phi^{+}}$ 有
        $X=\tfrac1{\sqrt2}\begin{psmallmatrix}1&0\\ 0&1\end{psmallmatrix}$，故
        $\det X=\tfrac12\neq0$：纠缠的——这是 \cref{thm:lec01-bell} 的一行式替代证明。
\end{itemize}
\end{example}

\begin{remark}[乘积态是例外，而非常态]\label{rem:lec01-generic-ent}
乘积态分解为相互独立的单量子比特态，故它不携带量子比特之间的任何关联：一个 $n$ 量子比特乘积态
$\ket{\psi_1}\otimes\cdots\otimes\ket{\psi_n}$ 由布洛赫球 (Bloch sphere) 上的 $n$ 个点确定，
仅 $2n$ 个实参数（见图）。一个一般的归一化 $n$ 量子比特态有 $2^{n}$ 个复振幅——除去全局相位后为
$2\cdot2^{n}-2$ 个实参数，呈指数多。对 $n=2$，这是一个 $4$-参数的乘积态族嵌于一个 $6$-参数的态空间
之内：乘积态构成一个低维（零测度）曲面，故\emph{几乎每一个}多量子比特态都是纠缠的。纠缠是常态；
乘积态是特殊的、“可经典描述的”一隅。
\end{remark}

\begin{center}
\begin{tikzpicture}[scale=0.9]
  \draw (0,0) circle (0.7);
  \draw[dashed] (-0.7,0) arc (180:360:0.7 and 0.24);
  \draw (0.7,0) arc (0:180:0.7 and 0.24);
  \draw[->,thick,red] (0,0) -- (0.32,0.5);
  \node at (0,-1.05) {$\ket{\psi_1}$};
  \node at (1.35,0) {$\otimes$};
  \begin{scope}[xshift=2.7cm]
    \draw (0,0) circle (0.7);
    \draw[dashed] (-0.7,0) arc (180:360:0.7 and 0.24);
    \draw (0.7,0) arc (0:180:0.7 and 0.24);
    \draw[->,thick,red] (0,0) -- (-0.45,0.35);
    \node at (0,-1.05) {$\ket{\psi_2}$};
  \end{scope}
  \node at (4.15,0) {$\otimes\ \cdots$};
  \node[align=left] at (7.3,0) {\small 一个乘积态：\\ \small $n$ 个相互独立的布洛赫箭头，\\ \small 没有共享的关联。};
\end{tikzpicture}
\end{center}

\subsection{贝尔态与纠缠}
\label{sec:lec01-bell-sub}

\begin{definition}[贝尔态]\label{def:lec01-bell}
\emph{贝尔态}（或 EPR 对）是两量子比特态
\begin{equation}
  \ket{\Phi^{+}}=\tfrac{1}{\sqrt2}\bigl(\ket{00}+\ket{11}\bigr).
\end{equation}
\end{definition}

\begin{theorem}[贝尔态是纠缠的]\label{thm:lec01-bell}
不存在一对单量子比特态 $\ket{\alpha},\ket{\beta}\in\C^{2}$ 使得
$\ket{\Phi^{+}}=\ket{\alpha}\otimes\ket{\beta}$。
\end{theorem}

\begin{proof}
反设存在这样的分解。写 $\ket{\alpha}=a_0\ket{0}+a_1\ket{1}$、
$\ket{\beta}=b_0\ket{0}+b_1\ket{1}$。则
\[
  \ket{\alpha}\otimes\ket{\beta}
    =a_0 b_0\ket{00}+a_0 b_1\ket{01}+a_1 b_0\ket{10}+a_1 b_1\ket{11}.
\]
与 $\ket{\Phi^{+}}=\tfrac{1}{\sqrt2}\ket{00}+0\cdot\ket{01}+0\cdot\ket{10}+\tfrac{1}{\sqrt2}\ket{11}$
匹配系数要求
\[
  a_0 b_0=\tfrac{1}{\sqrt2},\qquad a_0 b_1=0,\qquad a_1 b_0=0,\qquad a_1 b_1=\tfrac{1}{\sqrt2}.
\]
由 $a_0 b_0=\tfrac{1}{\sqrt2}\neq0$ 得 $a_0\neq0$ 且 $b_0\neq0$。于是
$a_0 b_1=0$ 迫使 $b_1=0$，而 $a_1 b_0=0$ 迫使 $a_1=0$；因此
$a_1 b_1=0\neq\tfrac{1}{\sqrt2}$，矛盾。故不存在任何分解，$\ket{\Phi^{+}}$ 是纠缠的。
\end{proof}

\begin{remark}[纠缠的物理意义]
$\ket{\Phi^{+}}$ 的关联没有可分的描述。在计算基中测量 (measurement) 第一个量子比特，各以
$\tfrac12$ 的概率得到 $0$ 或 $1$；无论结果如何，第二个量子比特的测量后态都被强制为\emph{相同}的
值（$0\mapsto\ket{0}$，$1\mapsto\ket{1}$）。即便两个量子比特在空间上相距甚远，这两个结果也完美
关联。这种无法表达为任何独立单量子比特态的经典混合的关联，正是隐形传态 (teleportation)、
超密编码 (superdense coding) 以及大部分量子优势背后的资源。
\end{remark}

\subsection{非定域性与 CHSH 不等式}
\label{sec:lec01-chsh}

贝尔对的关联不仅仅是强；它们\emph{比任何遵从定域实在论的经典机制所能产生的都要强}。这正是
贝尔定理的内容，我们将以 Clauser、Horne、Shimony 和 Holt（CHSH）给出的锐利形式来展开它。
值得注意的一点是，一个哲学假设——测量结果是预先确定的，且不受远处选择的影响——竟作出了一个
\emph{定量的}、可检验的预言，而量子力学违反了它。

两个被分开的参与者 Penelope 与 Odysseus 各自接收一对共享对中的一个量子比特。每人持有一台
具有两种设置的测量装置。Penelope 的两种可能测量返回值 $A_1,A_2\in\{+1,-1\}$；Odysseus 的
返回 $B_1,B_2\in\{+1,-1\}$。每一轮中，各方自由地选取一种设置，故 $A_1,A_2$ 中只有一个、
$B_1,B_2$ 中只有一个被实际观测到。考虑如下组合
\begin{equation}\label{eq:lec01-chsh-S}
  S \;=\; A_1B_1-A_1B_2+A_2B_1+A_2B_2 \;=\; A_1(B_1-B_2)+A_2(B_1+B_2).
\end{equation}

\begin{theorem}[CHSH 不等式]\label{thm:lec01-chsh}
假设四个量 $A_1,A_2,B_1,B_2\in\{+1,-1\}$ 在每一轮中都具有确定的值（\emph{实在论}），且
Penelope 的值不依赖于 Odysseus 的设置，反之亦然（\emph{定域性}）。则 $S$ 的期望满足
\[
  \bigl|\,\E[S]\,\bigr|\;=\;\bigl|\,\E[A_1B_1]-\E[A_1B_2]+\E[A_2B_1]+\E[A_2B_2]\,\bigr|\;\le\;2.
\]
\end{theorem}

\begin{proof}
在实在论下，四个值在每一轮中共存，故 \eqref{eq:lec01-chsh-S} 可逐项成立。由于
$B_1,B_2\in\{+1,-1\}$，$B_1-B_2$ 与 $B_1+B_2$ 中恰有一个为 $0$，另一个为 $\pm2$。因此在每一轮中，
两个乘积 $A_1(B_1-B_2)$、$A_2(B_1+B_2)$ 中有一个为零，另一个等于 $\pm2$，故 $S\in\{+2,-2\}$ 是
确定性的。一个始终为 $\pm2$ 的量，其平均落在 $[-2,2]$ 中，而期望的线性性把 $\E[S]$ 拆成四个
关联子。
\end{proof}

量子力学预言一个更大的值。在布洛赫球的 $xz$-平面中“沿角度 $\theta$”测量一个量子比特，意指测量
以 $\ket{\theta}=\cos\tfrac{\theta}{2}\ket{0}+\sin\tfrac{\theta}{2}\ket{1}$（结果 $+1$）及其正交
伙伴（结果 $-1$）为本征态的可观测量；在操作上，先作用把 $\ket{\theta}\mapsto\ket{0}$ 的旋转
(rotation) $R_y(-\theta)$，再在计算基中测量。

\begin{proposition}[贝尔对的量子关联子]\label{prop:lec01-chsh-corr}
若 Penelope 沿角度 $\alpha$ 测量她那一半 $\ket{\Phi^{+}}=\tfrac{1}{\sqrt2}(\ket{00}+\ket{11})$，
而 Odysseus 沿角度 $\beta$ 测量他那一半（采用上面的 $\pm1$ 标记），则其结果的关联为
\begin{equation}\label{eq:lec01-chsh-cosab}
  \E[AB]\;=\;\cos(\alpha-\beta).
\end{equation}
\end{proposition}

\begin{proof}
把 $R_y(-\alpha)\otimes R_y(-\beta)$ 作用于 $\ket{\Phi^{+}}$，并利用
$R_y(-\theta)\ket{0}=\cos\tfrac{\theta}{2}\ket{0}-\sin\tfrac{\theta}{2}\ket{1}$、
$R_y(-\theta)\ket{1}=\sin\tfrac{\theta}{2}\ket{0}+\cos\tfrac{\theta}{2}\ket{1}$，
展开并整理各项。半角加法恒等式
$\cos\tfrac{\alpha}{2}\cos\tfrac{\beta}{2}+\sin\tfrac{\alpha}{2}\sin\tfrac{\beta}{2}
=\cos\tfrac{\alpha-\beta}{2}$ 与
$\sin\tfrac{\alpha}{2}\cos\tfrac{\beta}{2}-\cos\tfrac{\alpha}{2}\sin\tfrac{\beta}{2}
=\sin\tfrac{\alpha-\beta}{2}$ 把该态化简为
\[
  \tfrac{1}{\sqrt2}\Bigl(\cos\tfrac{\alpha-\beta}{2}\,\ket{00}
      +\sin\tfrac{\alpha-\beta}{2}\,\ket{01}
      -\sin\tfrac{\alpha-\beta}{2}\,\ket{10}
      +\cos\tfrac{\alpha-\beta}{2}\,\ket{11}\Bigr).
\]
结果 $00,11$（乘积 $+1$）各以概率 $\tfrac12\cos^{2}\tfrac{\alpha-\beta}{2}$ 出现，而
$01,10$（乘积 $-1$）各以概率 $\tfrac12\sin^{2}\tfrac{\alpha-\beta}{2}$ 出现。因此由倍角恒等式，
$\E[AB]=\cos^{2}\tfrac{\alpha-\beta}{2}-\sin^{2}\tfrac{\alpha-\beta}{2}=\cos(\alpha-\beta)$。
\end{proof}

\begin{corollary}[Tsirelson 点的违反]\label{cor:lec01-tsirelson}
在 Penelope 取设置 $\alpha_1=0,\ \alpha_2=\tfrac{\pi}{2}$、Odysseus 取
$\beta_1=\tfrac{\pi}{4},\ \beta_2=\tfrac{3\pi}{4}$ 时，量子预言为
\[
  \E[S]=\cos(\alpha_1-\beta_1)-\cos(\alpha_1-\beta_2)+\cos(\alpha_2-\beta_1)+\cos(\alpha_2-\beta_2)
       =4\cos\tfrac{\pi}{4}=2\sqrt2\approx2.828,
\]
它超过了 \cref{thm:lec01-chsh} 的经典界 $2$。
\end{corollary}

\begin{proof}
四个角度差各为 $\mp\tfrac{\pi}{4}$：具体地，
$\alpha_1-\beta_1=-\tfrac{\pi}{4}$、$\alpha_1-\beta_2=-\tfrac{3\pi}{4}$、
$\alpha_2-\beta_1=\tfrac{\pi}{4}$、$\alpha_2-\beta_2=-\tfrac{\pi}{4}$。由于
$\cos(\pm\tfrac{\pi}{4})=\tfrac{1}{\sqrt2}$ 且 $\cos(\tfrac{3\pi}{4})=-\tfrac{1}{\sqrt2}$，
\[
  \E[S]=\tfrac{1}{\sqrt2}-\bigl(-\tfrac{1}{\sqrt2}\bigr)+\tfrac{1}{\sqrt2}+\tfrac{1}{\sqrt2}
       =\tfrac{4}{\sqrt2}=2\sqrt2. \qedhere
\]
\end{proof}

\begin{remark}[违反迫使我们放弃什么]
$2\sqrt2$ 是任何量子态与测量所能达到的最大值（Tsirelson 界），并已被实验证实——最著名的是
2022 年诺贝尔物理学奖所表彰的无漏洞检验。由于 \cref{thm:lec01-chsh} 仅由实在论加定域性推出，
它的实验违反意味着这两个假设中至少有一个不成立：测量前的值 $A_1,A_2,B_1,B_2$ 不可能全部作为
定域的“隐变量”同时存在。$\ket{\Phi^{+}}$ 的关联是真正非经典的，然而——正如我们接下来在
\cref{sec:lec01-tele} 中所示——它们仍然不能用来传递信号，故与相对论并无冲突。
\end{remark}
\section{门、通用门组与线路}
\label{sec:lec01-3}

量子力学的公设 (postulate) 告诉我们，孤立系统的态随时间通过一个线性、保内积的映射
演化。在有限维希尔伯特空间上，这样的映射恰好就是一个酉算符 (unitary operator)，在量子
计算中我们称之为门 (gate)。

\begin{definition}[酉算符 / 量子门]\label{def:lec01-unitary}
设 $\Hilb\cong\C^{d}$ 为有限维希尔伯特空间。线性算符
$U:\Hilb\to\Hilb$ 称为\emph{酉的}，若
\[
  U\dg U = U U\dg = \id,
\]
其中 $U\dg$ 是 $U$ 的共轭转置（伴随）。$n$ 个量子比特 (qubit) 上的\emph{量子门}是
$\Hilb=(\C^{2})^{\otimes n}\cong\C^{2^{n}}$ 上的酉算符。
\end{definition}

等价地，$U$ 是酉的当且仅当其列（或行）构成 $\C^{d}$ 的一组标准正交基，当且仅当
$U^{-1}=U\dg$。特别地，每个门都是可逆的，因此量子计算是可逆的 (reversible)。

\begin{theorem}[门保持内积与范数]\label{thm:lec01-unitary-preserves}
设 $U$ 为 $\Hilb$ 上的线性算符。以下命题等价。
\begin{enumerate}[label=\textup{(\roman*)}]
  \item $U$ 是酉的，即 $U\dg U=\id$。
  \item $U$ 保持内积 (inner product)：对所有 $\ket{\psi},\ket{\phi}\in\Hilb$ 有
        $\braket{U\psi|U\phi}=\braket{\psi|\phi}$。
  \item $U$ 保持范数：对所有 $\ket{\psi}\in\Hilb$ 有 $\norm{U\psi}=\norm{\psi}$。
\end{enumerate}
\end{theorem}

\begin{proof}
$(\mathrm{i})\Rightarrow(\mathrm{ii})$。利用
$\braket{U\psi|U\phi}=\bra{\psi}U\dg U\ket{\phi}$ 以及 $U\dg U=\id$，得到
$\braket{U\psi|U\phi}=\bra{\psi}\id\ket{\phi}=\braket{\psi|\phi}$。

$(\mathrm{ii})\Rightarrow(\mathrm{iii})$。取 $\ket{\phi}=\ket{\psi}$：则
$\norm{U\psi}^{2}=\braket{U\psi|U\psi}=\braket{\psi|\psi}=\norm{\psi}^{2}$，再
取非负平方根即得 $\norm{U\psi}=\norm{\psi}$。

$(\mathrm{iii})\Rightarrow(\mathrm{i})$。范数保持给出
对所有 $\ket{\psi}$ 有 $\bra{\psi}U\dg U\ket{\psi}=\bra{\psi}\id\ket{\psi}$，故
自伴算符 $M\defeq U\dg U-\id$ 对每个 $\ket{\psi}$ 满足 $\bra{\psi}M\ket{\psi}=0$。
对自伴算符而言，二次型决定了算符（通过极化恒等式，或取 $\ket{\psi}$ 为其
本征向量 (eigenvector)），因此 $M=0$，即 $U\dg U=\id$。在有限维空间上左
逆就是双边逆，故也有 $UU\dg=\id$，于是 $U$ 是酉的。
\end{proof}

从物理上看，\cref{thm:lec01-unitary-preserves} 说明门将归一化的态映为归一化的
态，且绝不改变两个态之间的``夹角''；正是这种内积的刚性，使得\cref{sec:lec01-4}
中的克隆成为不可能。

\subsection{单量子比特门}

\begin{definition}[泡利门、Hadamard 门与相位门]\label{def:lec01-gates}
$\C^{2}$ 上的\emph{泡利门 (Pauli gate)}与恒等算符为
\[
  I=\begin{pmatrix}1&0\\0&1\end{pmatrix},\quad
  X=\begin{pmatrix}0&1\\1&0\end{pmatrix},\quad
  Y=\begin{pmatrix}0&-\ii\\\ii&0\end{pmatrix},\quad
  Z=\begin{pmatrix}1&0\\0&-1\end{pmatrix}.
\]
\emph{Hadamard 门}与\emph{相位门} $S$ 和 $T$ 为
\[
  H=\frac{1}{\sqrt{2}}\begin{pmatrix}1&1\\1&-1\end{pmatrix},\qquad
  S=\begin{pmatrix}1&0\\0&\ii\end{pmatrix},\qquad
  T=\begin{pmatrix}1&0\\0&\ee^{\ii\pi/4}\end{pmatrix}.
\]
\end{definition}

它们每个都是酉的：直接验证可知 $X\dg X=Y\dg Y=Z\dg Z=I$，$H\dg H=I$，以及
$S,T$ 为对角矩阵且对角元模长为 $1$。在计算基下，
\[
  X\ket{0}=\ket{1},\quad X\ket{1}=\ket{0},\qquad
  Z\ket{0}=\ket{0},\quad Z\ket{1}=-\ket{1},
\]
因此 $X$ 是\emph{比特翻转}（量子 \textsc{not} 门），而 $Z$ 是\emph{相位翻转}：
它固定 $\ket{0}$ 并将 $\ket{1}$ 乘以 $-1$。Hadamard 门制造叠加 (superposition)，
\[
  H\ket{0}=\ket{+}=\tfrac{1}{\sqrt{2}}(\ket{0}+\ket{1}),\qquad
  H\ket{1}=\ket{-}=\tfrac{1}{\sqrt{2}}(\ket{0}-\ket{1}),
\]
且是自逆的，$H^{2}=I$。相位门保持 $\ket{0}$ 不变并旋转 $\ket{1}$ 的相位：
$S$ 将其乘以 $\ii$，$T$ 将其乘以 $\ee^{\ii\pi/4}$。

四个算符 $\{I,X,Y,Z\}$ 构成 $2\times2$ 矩阵的（复）向量空间的一组基，因此每个
单量子比特门都是泡利矩阵的复线性组合。从几何上看，单量子比特纯态对应于
布洛赫球 (Bloch sphere) 上的一点；泡利门作为绕 $x$、$y$、$z$ 轴旋转 $\pi$ 的
旋转 (rotation) 来作用（相差一个全局相位）。这些半圈旋转是一个单参数族的特例，
我们现在从泡利代数中将其推导出来。对于单位轴 $\hat n\in\R^{3}$，记
$\hat n\cdot\vec\sigma\defeq n_xX+n_yY+n_zZ$。

\begin{proposition}[旋转门及其泡利生成元]\label{prop:lec01-rot}
泡利乘积恒等式
$(\vec a\cdot\vec\sigma)(\vec b\cdot\vec\sigma)=(\vec a\cdot\vec b)\,I+\ii(\vec a\times\vec b)\cdot\vec\sigma$
给出 $(\hat n\cdot\vec\sigma)^{2}=I$，因此绕 $\hat n$ 转过角度 $\theta$ 的
\emph{旋转门}具有闭式
\[
  R_{\hat n}(\theta)\defeq\exp\!\Big(-\tfrac{\ii\theta}{2}\,\hat n\cdot\vec\sigma\Big)
  =\cos\tfrac{\theta}{2}\,I-\ii\sin\tfrac{\theta}{2}\,(\hat n\cdot\vec\sigma).
\]
取 $\hat n=\hat x,\hat y,\hat z$ 得到
\[
  R_x(\theta)=\begin{pmatrix}\cos\tfrac{\theta}{2}&-\ii\sin\tfrac{\theta}{2}\\[2pt]
      -\ii\sin\tfrac{\theta}{2}&\cos\tfrac{\theta}{2}\end{pmatrix},\quad
  R_y(\theta)=\begin{pmatrix}\cos\tfrac{\theta}{2}&-\sin\tfrac{\theta}{2}\\[2pt]
      \sin\tfrac{\theta}{2}&\cos\tfrac{\theta}{2}\end{pmatrix},\quad
  R_z(\theta)=\begin{pmatrix}\ee^{-\ii\theta/2}&0\\0&\ee^{\ii\theta/2}\end{pmatrix},
\]
即 \eqref{eq:lec01-rotations} 中用于态制备 (state preparation) 的门。泡利门是
半圈旋转：$X=\ii R_x(\pi)$，$Y=\ii R_y(\pi)$，$Z=\ii R_z(\pi)$。
\end{proposition}

\begin{proof}
由于 $(\hat n\cdot\vec\sigma)^{2}=I$，在级数
$\exp(-\ii\alpha\,\hat n\cdot\vec\sigma)=\sum_{k\ge0}\tfrac{(-\ii\alpha)^{k}}{k!}(\hat n\cdot\vec\sigma)^{k}$
（其中 $\alpha=\theta/2$）中偶次幂约化为 $I$，奇次幂约化为
$\hat n\cdot\vec\sigma$，因此它们重新求和为 $\cos\alpha\,I-\ii\sin\alpha\,(\hat n\cdot\vec\sigma)$。
代入 $X,Y,Z$ 给出三个矩阵；例如
$R_z(\theta)=\cos\tfrac{\theta}{2}I-\ii\sin\tfrac{\theta}{2}Z=\diag(\ee^{-\ii\theta/2},\ee^{\ii\theta/2})$。
最后 $R_z(\pi)=\diag(\ee^{-\ii\pi/2},\ee^{\ii\pi/2})=-\ii Z$，故 $Z=\ii R_z(\pi)$，
$X,Y$ 同理。
\end{proof}

\begin{proposition}[$R_{\hat n}(\theta)$ 旋转布洛赫向量]\label{prop:lec01-rodrigues}
将单量子比特态写为 $\rho=\tfrac12(I+\vec r\cdot\vec\sigma)$。用 $R_{\hat n}(\theta)$
共轭将布洛赫向量变为其绕 $\hat n$ 转过角度 $\theta$ 的旋转：
\[
  R_{\hat n}(\theta)\,\rho\,R_{\hat n}(\theta)\dg=\tfrac12\big(I+\vec r\,'\cdot\vec\sigma\big),
  \quad
  \vec r\,'=\cos\theta\,\vec r+\sin\theta\,(\hat n\times\vec r)+(1-\cos\theta)(\hat n\cdot\vec r)\,\hat n .
\]
\end{proposition}

\begin{proof}
由于 $R\,I\,R\dg=I$，只有 $\vec r\cdot\vec\sigma$ 部分发生移动，故只需对一般的
$\vec m$ 共轭 $M\defeq\vec m\cdot\vec\sigma$。令 $c=\cos\tfrac{\theta}{2}$，
$s=\sin\tfrac{\theta}{2}$，$N\defeq\hat n\cdot\vec\sigma$，于是 $R=cI-\ii sN$ 且
$R\dg=cI+\ii sN$；则
\[
  RMR\dg=c^{2}M+\ii cs\,[M,N]+s^{2}\,NMN .
\]
乘积恒等式给出 $[M,N]=2\ii(\vec m\times\hat n)\cdot\vec\sigma$，经短暂计算后得
$NMN=\big(2(\hat n\cdot\vec m)\,\hat n-\vec m\big)\cdot\vec\sigma$。
因此 $RMR\dg=\vec m\,'\cdot\vec\sigma$，其中
\[
  \vec m\,'=c^{2}\vec m+2cs\,(\hat n\times\vec m)+s^{2}\big(2(\hat n\cdot\vec m)\,\hat n-\vec m\big),
\]
而倍角恒等式 $c^{2}-s^{2}=\cos\theta$、$2cs=\sin\theta$、
$2s^{2}=1-\cos\theta$ 将其收缩为罗德里格斯 (Rodrigues) 旋转公式。令
$\vec m=\vec r$ 即证得结论。
\end{proof}

\begin{remark}[半角与 $SU(2)\to SO(3)$ 双重覆盖]\label{rem:lec01-doublecover}
态携带半角 $\theta/2$，而布洛赫向量转过整个 $\theta$。于是
$R_{\hat n}(2\pi)=\cos\pi\,I-\ii\sin\pi\,(\hat n\cdot\vec\sigma)=-I$：$2\pi$ 的
布洛赫旋转使向量复原，却给态留下一个全局相位 $-1$，只有 $4\pi$ 的旋转才恢复 $I$。
等价地，$R_{\hat n}(\theta)\mapsto\mathcal R(\hat n,\theta)$ 是一个二对一同态
$SU(2)\to SO(3)$，其中 $\pm R_{\hat n}(\theta)$ 诱导相同的（可观测的）旋转——这
正是布洛赫图像能忠实地将单量子比特门表示为普通三维旋转的原因。
\end{remark}

\begin{lemma}[Hadamard 共轭]\label{lem:lec01-conj}
$H X H = Z$ 且 $H Z H = X$。
\end{lemma}

\begin{proof}
回忆 $H=H\dg$ 且 $H^{2}=I$。通过矩阵乘法计算 $HXH$。首先
\[
  XH=\begin{pmatrix}0&1\\1&0\end{pmatrix}
     \frac{1}{\sqrt2}\begin{pmatrix}1&1\\1&-1\end{pmatrix}
    =\frac{1}{\sqrt2}\begin{pmatrix}1&-1\\1&1\end{pmatrix},
\]
然后
\[
  HXH=\frac{1}{\sqrt2}\begin{pmatrix}1&1\\1&-1\end{pmatrix}\cdot
      \frac{1}{\sqrt2}\begin{pmatrix}1&-1\\1&1\end{pmatrix}
     =\frac{1}{2}\begin{pmatrix}2&0\\0&-2\end{pmatrix}
     =\begin{pmatrix}1&0\\0&-1\end{pmatrix}=Z.
\]
对于第二个恒等式，用 $H$ 对两边共轭第一个恒等式并利用 $H^{2}=I$：由
$HXH=Z$ 得 $H(HXH)H=HZH$，即 $(H^{2})X(H^{2})=HZH$，故 $X=HZH$。
\end{proof}

\Cref{lem:lec01-conj} 精确地表明 $H$ 是比特翻转图像与相位翻转图像之间的
``桥梁''：计算基中的相位翻转，在 Hadamard（$\ket{+},\ket{-}$）基中看起来像比特
翻转，反之亦然。另有两个恒等式记录于此以备参考。

\begin{proposition}\label{prop:lec01-yxz}
$Y=\ii XZ$ 且 $S=T^{2}$。
\end{proposition}

\begin{proof}
直接计算：
\[
  XZ=\begin{pmatrix}0&1\\1&0\end{pmatrix}
     \begin{pmatrix}1&0\\0&-1\end{pmatrix}
    =\begin{pmatrix}0&-1\\1&0\end{pmatrix},\qquad
  \ii XZ=\begin{pmatrix}0&-\ii\\\ii&0\end{pmatrix}=Y.
\]
对于相位门，$T$ 为对角矩阵，故 $T^{2}=\diag(1,\ee^{\ii\pi/4})^{2}=\diag(1,\ee^{\ii\pi/2})=\diag(1,\ii)=S$。
\end{proof}

\begin{remark}[$T$ 是非 Clifford 的]
\emph{Clifford 群}由 $\{H,S,\CNOT\}$ 生成（等价地，即在共轭下将泡利群映到自身的
那些酉算符）。上面的恒等式表明 $X,Y,Z$ 和 $S=T^{2}$ 落在这个群中。而 $T$ 门\emph{不}
在其中：用 $T$ 共轭一个泡利算符产生的是泡利算符的叠加，而非另一个泡利算符。这一点
很重要，因为 Gottesman--Knill 定理指出 Clifford 线路可被经典高效模拟；$T$ 门正是
将门组提升为量子计算通用性的成分。
\end{remark}

\subsection{多量子比特门}

要作用于不止一个量子比特，我们采用张量积。作用于量子比特~$0$ 上的门 $A$ 和作用于
量子比特~$1$ 上的门 $B$ 以 $A\otimes B$ 作用于联合态，其中
$(A\otimes B)(\ket{u}\otimes\ket{v})=A\ket{u}\otimes B\ket{v}$。最重要的真正
双量子比特门是受控 \textsc{not}。

\begin{definition}[$\CNOT$ 与 Toffoli]\label{def:lec01-cnot}
\emph{受控 \textsc{not}} 门 $\CNOT$ 作用于两个量子比特（控制位 $a$，目标位 $b$），
\[
  \CNOT\ket{a,b}=\ket{a,\,a\oplus b},\qquad a,b\in\{0,1\},
\]
其中 $\oplus$ 是模 $2$ 加法。在有序计算基
$\ket{00},\ket{01},\ket{10},\ket{11}$ 下，其矩阵是 $4\times4$ 置换
酉矩阵
\[
  \CNOT=\begin{pmatrix}
    1&0&0&0\\ 0&1&0&0\\ 0&0&0&1\\ 0&0&1&0
  \end{pmatrix}.
\]
\emph{Toffoli}（或 \textsc{ccnot}）门作用于三个量子比特，
$\mathrm{CCNOT}\ket{a,b,c}=\ket{a,\,b,\,(a\wedge b)\oplus c}$，当且仅当两个控制位
$a,b$ 都等于 $1$ 时翻转目标位 $c$。
\end{definition}

由于 $\CNOT$ 置换四个基向量，它是一个置换矩阵，因而是酉的；Toffoli 门同理，它交换
$\ket{110}\leftrightarrow\ket{111}$。两者都是自逆的。Toffoli 门对经典可逆计算是
通用的：借助辅助比特，它可以实现 \textsc{and}、\textsc{not} 和扇出，因此任何布尔
线路都可被改造为可逆的。

\begin{remark}[为何异或保持谕示可逆]
$f:\{0,1\}^{n}\to\{0,1\}$ 的经典谕示 (oracle) 在量子上实现为
$\Uf\ket{x}\ket{b}=\ket{x}\ket{b\oplus f(x)}$。用 $\oplus$ 写入答案而非用 $f(x)$
覆写 $b$，正是使 $\Uf$ 成为基态的置换、从而是酉的且可逆的原因：施加 $\Uf$ 两次
返回 $\ket{b\oplus f(x)\oplus f(x)}=\ket{b}$，故 $\Uf^{2}=\id$。单纯的覆写
$b\mapsto f(x)$ 一般不是单射，无法成为酉的。$\CNOT$ 规则 $b\mapsto a\oplus b$ 正是
这一构造的单比特实例。
\end{remark}

\subsection{阅读量子线路}

\emph{量子线路 (circuit)}是一个图示，其中每条水平导线是一个量子比特，门绘制在
导线上。有三条约定支配着如何阅读线路。
\begin{enumerate}[label=(\arabic*)]
  \item \textbf{时间从左向右流动。}输入态位于最左端；绘制得更靠左的门更早
        被施加。
  \item \textbf{控制位。}在受控门上，实心点 $\bullet$ 标记一个在量子比特为
        $\ket{1}$ 时激活的控制位；空心点 $\circ$ 标记一个 $0$-控制位，它在量子
        比特为 $\ket{0}$ 时激活。
  \item \textbf{矩阵顺序是反过来的。}若线路先施加 $U_{1}$，再施加
        $U_{2}$，则整体酉算符是\emph{逆序}乘积 $U=U_{2}\,U_{1}$，
        因为算符作用于其右侧的右矢。对于长度为 $m$ 的序列，
        $U=U_{m}U_{m-1}\cdots U_{1}$。
\end{enumerate}

控制位为 $a$ 的受控 $U$ 门可以简洁地写为
$c\text{-}U\,\ket{a}\ket{b}=\ket{a}\,U^{a}\ket{b}$，即 $U$ 以 $a\in\{0,1\}$ 为
幂次作用于目标位（当 $a=0$ 时什么也不发生）。\textsc{swap} 门交换两个量子比特，
$\mathrm{SWAP}\ket{a,b}=\ket{b,a}$；它可由三个 $\CNOT$ 构建。

\begin{definition}[线路深度]\label{def:lec01-depth}
将一条线路的门分组为\emph{层}，其中一层是作用于互不相交量子比特、可以同时执行的
最大门集合。线路的\emph{深度 (circuit depth)}是层数，等价地是穿过图示的任一
输入到输出路径上门数的最大值。
\end{definition}

作用于互不相交量子比特集合的门作为张量因子彼此对易并并行运行，因此是深度（而非
总门数 (gate count)）衡量了墙钟运行时间，而保持低深度对于在含噪硬件上运行至关重要。

\begin{example}[将小线路的酉算符写为逆序乘积]
考虑一个量子比特，从左到右依次为一个 $H$ 门、一个 $S$ 门，再一个 $H$ 门。整体
酉算符是逆序乘积
\[
  U = H\,S\,H
    = \frac{1}{2}\begin{pmatrix}1&1\\1&-1\end{pmatrix}
      \begin{pmatrix}1&0\\0&\ii\end{pmatrix}
      \begin{pmatrix}1&1\\1&-1\end{pmatrix}
    = \frac{1}{2}\begin{pmatrix}1+\ii&1-\ii\\1-\ii&1+\ii\end{pmatrix}.
\]
验证 $U\dg U=I$，确认该复合算符仍是一个门。
\end{example}

\subsection{构建贝尔态}

我们现在将单量子比特门和双量子比特门组装成一条产生纠缠 (entanglement) 的线路。
将两个量子比特制备为 $\ket{00}$，对量子比特~$0$ 施加 $H$，然后以量子比特~$0$ 为
控制位、量子比特~$1$ 为目标位施加 $\CNOT$。

\begin{theorem}[贝尔态制备]\label{thm:lec01-bellprep}
以量子比特~$0$ 为控制位、量子比特~$1$ 为目标位，
\[
  \CNOT\,(H\otimes I)\,\ket{00}
    = \tfrac{1}{\sqrt2}\big(\ket{00}+\ket{11}\big)=\ket{\Phi^{+}}.
\]
\end{theorem}

\begin{proof}
首先施加 $H\otimes I$。由于 $H\ket{0}=\tfrac{1}{\sqrt2}(\ket{0}+\ket{1})$ 且
$I\ket{0}=\ket{0}$，
\[
  (H\otimes I)\ket{00}
    = \tfrac{1}{\sqrt2}(\ket{0}+\ket{1})\otimes\ket{0}
    = \tfrac{1}{\sqrt2}\big(\ket{00}+\ket{10}\big).
\]
现在利用 $\CNOT\ket{a,b}=\ket{a,a\oplus b}$ 逐项施加 $\CNOT$。在 $\ket{00}$ 上
控制位为 $0$，故态不变。在 $\ket{10}$ 上控制位为 $1$，故目标位翻转：
$\ket{10}\mapsto\ket{11}$。因此
\[
  \CNOT\,\tfrac{1}{\sqrt2}\big(\ket{00}+\ket{10}\big)
    = \tfrac{1}{\sqrt2}\big(\ket{00}+\ket{11}\big)=\ket{\Phi^{+}}. \qedhere
\]
\end{proof}

根据 \cref{thm:lec01-bell}，态 $\ket{\Phi^{+}}$ 是纠缠的，因此这条两个门的线路
将乘积输入变为了一个最大纠缠的贝尔对 (Bell pair)。

\subsection{通用门组与编译}

\begin{definition}[通用门组]\label{def:lec01-universal}
一个有限门集合 $\mathcal{G}$ 称为\emph{通用的}，若将 $\mathcal{G}$ 中的门施加于
任意量子比特子集所生成的酉算符，在所有 $n$ 量子比特酉算符构成的群中稠密（相差
一个全局相位）：对于每个目标酉算符 $U$ 和每个 $\varepsilon>0$，都存在一条基于
$\mathcal{G}$ 的有限线路 $V$，使得 $\norm{U-V}\le\varepsilon$。
\end{definition}

两个事实支撑着这个定义。第一，任何 $n$ 量子比特酉算符都可以\emph{精确地}写为
单量子比特门和双量子比特门的乘积，$U=U_{m}U_{m-1}\cdots U_{1}$，尽管 $m$ 可能是
$n$ 的指数。第二，一个固定的\emph{有限}集合就已足以逼近每个门：集合
$\mathcal{G}=\{H,\,T,\,\CNOT\}$ 是通用的。我们满足于逼近，是因为一个离散集合
无法精确命中连续统中的酉算符，但它可以逼近到任意 $\varepsilon$ 之内。

\begin{remark}[Solovay--Kitaev 代价]
Solovay--Kitaev 定理量化了这种离散化的代价：任何单量子比特门都可以被一条由通用
集合中 $\bigo\!\big(\polylog(1/\varepsilon)\big)$ 个门构成的线路逼近到误差
$\varepsilon$，而一条 $m$ 个门的线路可以仅用
$\bigo\!\big(m\,\polylog(m/\varepsilon)\big)$ 个门重新编译到有限集合中，并由一个
高效的经典算法找到。因此开销仅在逆精度上呈多重对数增长，这正是将理想线路编译到
固定硬件门组切实可行的原因。
\end{remark}

\begin{remark}[Solovay--Kitaev 为何有效（概述）]\label{rem:lec01-sk-sketch}
多重对数界来自一个递归，它通过用越来越小的门来表示逼近的\emph{误差}从而使逼近
越来越精细。假设 $\mathcal{G}$ 对取逆封闭。
\emph{(i) 基网。}$\langle\mathcal{G}\rangle$ 的稠密性连同 $SU(2)$ 的紧性，保证了
用有界长度的词即可使\emph{每个}门达到一个固定精度 $\varepsilon_0$。
\emph{(ii) 对易子引理。}任何靠近恒等的 $\Delta$，$\norm{\Delta-\id}\le\delta$，
可分解为群对易子 $\Delta=VWV\dg W\dg$，误差为 $\bigo(\delta^{3/2})$，其中
$V,W$ 距离恒等仅为 $\bigo(\sqrt{\delta})$——写
$V=\ee^{-\ii A}$，$W=\ee^{-\ii B}$ 且 $A,B=\bigo(\sqrt{\delta})$，则有
$VWV\dg W\dg=\ee^{-[A,B]+\bigo(\delta^{3/2})}$ 且 $[A,B]=\bigo(\delta)$。
\emph{(iii) 递归。}给定精度为 $\varepsilon_{n-1}$ 的第 $(n{-}1)$ 层逼近
$U_{n-1}$，用该引理分解残差 $\Delta=U\,U_{n-1}\dg$，在第 $n{-}1$ 层逼近
$V,W$，并将对易子粘接到 $U_{n-1}$ 上：
\[
  \varepsilon_n=\bigo\!\big(\varepsilon_{n-1}^{3/2}\big),\qquad
  \ell_n=5\,\ell_{n-1}+\bigo(1).
\]
指数 $3/2$ 使误差呈双重指数下降，因此 $n=\bigo(\log\log(1/\varepsilon))$ 层即
足够；展开 $\ell_n=5^{n}\ell_0$ 便
给出 $\ell(\varepsilon)=\bigo\!\big((\log(1/\varepsilon))^{\log 5/\log(3/2)}\big)$，
指数 $\approx 3.97$。单量子比特旋转正是一般目标
$U\in SU(2)$，故该界直接适用于它（Dawson \& Nielsen,
\emph{The Solovay--Kitaev algorithm}, Quantum Inf.\ Comput.\ \textbf{6} (2006) 81）。
\end{remark}

在真实设备上运行一条抽象线路需要\emph{编译}：
(i) 将诸如 $H$、$T$、$\CNOT$ 的逻辑门重写为机器的原生门组；以及
(ii) 插入 \textsc{swap} 门以路由量子比特，因为硬件量子比特并非两两全连通。
整个过程中都要最小化门数和深度，以抑制噪声并最大化保真度。
\subsection{实例：可逆的 (reversible) 二进制加法器}
\label{sec:lec01-adder}

$n$ 个量子比特的计算基由整数 $0,1,\dots,2^{n}-1$ 标记：我们把字符串
$\ket{k_1\cdots k_n}$ 与其二进制值 $k=\sum_{l=1}^{n}k_l\,2^{\,n-l}$ 等同，于是
$\ket{10}=\ket{2}$、$\ket{11}=\ket{3}$，而三比特寄存器 (register) 遍历
$\ket{0},\dots,\ket{7}$。因此，一个置换这些基态的线路 \emph{计算了一个取整数值的
函数}，而最终的计算基测量则以普通数的形式返回答案。我们通过构造一个把两位数相加的
线路来将其具体化，这同时也验证了 \cref{def:lec01-cnot} 之后的论断：Toffoli 门
对经典可逆计算是通用的。这里 $\oplus$ 是模 $2$ 加法（即 \textsc{xor}，由 $\CNOT$
计算），而 $\wedge$ 是逻辑 \textsc{and}（由作用在一个 $\ket{0}$ 辅助比特 (ancilla)
上的 Toffoli 门计算）。

要把 $A_1A_0$ 与 $B_1B_0$ 相加，我们逐列进行，正如小学里的加法，只是在二进制下。
最右一列产生一个和位 $S_0$ 和一个进位 $C_1$；下一列把 $A_1$、$B_1$ 和进来的进位
$C_1$ 相加，产生 $S_1$ 和最终进位 $C_2$；结果是三位数 $C_2S_1S_0$（例如，$2+3$ 是
$10+11=101=5$）。在熟悉的竖式布局中，每个进位写在它所进入的那一列之上：
\[
\begin{array}{r@{\;}c@{\;}c@{\;}c}
    & C_2 & C_1 &     \\
    &     & A_1 & A_0 \\
  + &     & B_1 & B_0 \\ \hline
    & C_2 & S_1 & S_0
\end{array}
\qquad\text{例如}\qquad
\begin{array}{r@{\;}c@{\;}c@{\;}c}
    & 1 & 0 &   \\
    &   & 1 & 0 \\
  + &   & 1 & 1 \\ \hline
    & 1 & 0 & 1
\end{array}
\]

\paragraph{半加器。} 最右一列没有进来的进位，所以
\[
  S_0=A_0\oplus B_0,\qquad C_1=A_0\wedge B_0,
\]
真值表证实了这一点：
\begin{center}
\begin{tabular}{cc@{\qquad}cc}
\toprule
$A_0$ & $B_0$ & $C_1$ & $S_0$\\ \midrule
0 & 0 & 0 & 0\\
0 & 1 & 0 & 1\\
1 & 0 & 0 & 1\\
1 & 1 & 1 & 0\\ \bottomrule
\end{tabular}
\end{center}
和位是单个 $\CNOT$（控制 $A_0$，目标 $B_0$），而进位是一个 Toffoli 门（控制
$A_0,B_0$），它把 $A_0\wedge B_0$ 写到一个新的 $\ket{0}$ 辅助比特上。因为 $\CNOT$
会覆写 $B_0$，我们\emph{先}计算进位：
\begin{center}
\begin{quantikz}
\lstick{$\ket{A_0}$} & \ctrl{1} & \ctrl{1} & \qw \\
\lstick{$\ket{B_0}$} & \ctrl{1} & \targ{}  & \qw \\
\lstick{$\ket{0}$}   & \targ{}  & \qw      & \qw
\end{quantikz}
\end{center}
自上而下读取，输出线上分别保存着 $A_0$、和 $S_0=A_0\oplus B_0$，以及进位
$C_1=A_0\wedge B_0$。

\paragraph{全加器。} 下一列将三个比特 $C_1$、$A_1$、$B_1$ 相加：
\[
  S_1=C_1\oplus A_1\oplus B_1,\qquad
  C_2=(A_1\wedge B_1)\oplus(C_1\wedge A_1)\oplus(C_1\wedge B_1).
\]
这些规则由三个输入比特 $C_1,A_1,B_1$ 上的真值表证实：
\begin{center}
\begin{tabular}{ccc@{\qquad}cc}
\toprule
$C_1$ & $A_1$ & $B_1$ & $C_2$ & $S_1$\\ \midrule
0 & 0 & 0 & 0 & 0\\
0 & 0 & 1 & 0 & 1\\
0 & 1 & 0 & 0 & 1\\
0 & 1 & 1 & 1 & 0\\
1 & 0 & 0 & 0 & 1\\
1 & 0 & 1 & 1 & 0\\
1 & 1 & 0 & 1 & 0\\
1 & 1 & 1 & 1 & 1\\ \bottomrule
\end{tabular}
\end{center}
进位 $C_2$ 由三个 Toffoli 门累加到一个 $\ket{0}$ 辅助比特上，每个乘积项对应一个；
和 $S_1$ 则是两个以 $B_1$ 线为目标的 $\CNOT$（控制 $C_1$ 和 $A_1$）。和之前一样，
我们先计算进位，因为形成 $S_1$ 会覆写 $B_1$：
\begin{center}
\begin{quantikz}
\lstick{$\ket{C_1}$} & \qw      & \ctrl{3} & \ctrl{3} & \ctrl{2} & \qw      & \qw\\
\lstick{$\ket{A_1}$} & \ctrl{2} & \ctrl{2} & \qw      & \qw      & \ctrl{1} & \qw\\
\lstick{$\ket{B_1}$} & \ctrl{1} & \qw      & \ctrl{1} & \targ{}  & \targ{}  & \qw\\
\lstick{$\ket{0}$}   & \targ{}  & \targ{}  & \targ{}  & \qw      & \qw      & \qw
\end{quantikz}
\end{center}
现在输出线上分别保存着 $C_1$、$A_1$、和 $S_1$，以及辅助比特上的进位 $C_2$。

\paragraph{两位加法器。} 将这两个模块级联起来——半加器的进位输出 $C_1$ 馈入全加器——
就得到计算 $A_1A_0+B_1B_0$ 的完整线路：半加器从 $A_0,B_0$ 产生 $S_0$ 和 $C_1$，而
全加器消耗 $C_1,A_1,B_1$ 产生 $S_1$ 和 $C_2$。测量保存 $S_0$、$S_1$、$C_2$ 的三条
线即返回和的二进制位。例如，要计算 $3+3$，可对四个输入量子比特（初始化为 $\ket{0}$）
施加 $X$ 以置 $A_1A_0=B_1B_0=11$，并读出输出 $110_2=6$。完整线路画在
\cref{fig:lec01-adder2} 中。

\begin{figure}[htbp]
\centering
\begin{quantikz}
\lstick{$\ket{A_0}$} & \ctrl{1} & \ctrl{1} & \qw      & \qw      & \qw      & \qw      & \qw      & \qw \\
\lstick{$\ket{B_0}$} & \ctrl{1} & \targ{}  & \qw      & \qw      & \qw      & \qw      & \qw      & \qw \\
\lstick{$\ket{0}$}   & \targ{}  & \qw      & \qw      & \ctrl{3} & \ctrl{3} & \ctrl{1} & \qw      & \qw \\
\lstick{$\ket{A_1}$} & \qw      & \qw      & \ctrl{1} & \ctrl{2} & \qw      & \targ{}  & \ctrl{1} & \qw \\
\lstick{$\ket{B_1}$} & \qw      & \qw      & \ctrl{1} & \qw      & \ctrl{1} & \qw      & \targ{}  & \qw \\
\lstick{$\ket{0}$}   & \qw      & \qw      & \targ{}  & \targ{}  & \targ{}  & \qw      & \qw      & \qw
\end{quantikz}
\caption{两位行波加法器 $A_1A_0+B_1B_0\mapsto C_2S_1S_0$：半加器（前两个门）计算
$S_0,C_1$，级联的全加器计算 $S_1,C_2$（三个 Toffoli 门产生进位，接着两个 $\CNOT$
产生和）。输出线自上而下为：$A_0$、$S_0$、$C_1$、$C_1\oplus A_1$、$S_1$、$C_2$；
和为 $C_2S_1S_0$。改编自 Brody，图 7.12。}
\label{fig:lec01-adder2}
\end{figure}

\paragraph{运行加法器。} 在硬件上每个量子比特都从 $\ket{0}$ 开始，所以一个输入比特
通过在加法器之前放置一个前导 $X$（\textsc{not}）门来置为 $1$。
\Cref{fig:lec01-adder-runs} 展示了三种情形：$0+0$（无 $X$ 门）、$1+1=01+01$
（$A_0$ 和 $B_0$ 上各一个 $X$），以及 $1+2=01+10$（$A_0$ 和 $B_1$ 上各一个 $X$）；
测量 $S_0,S_1,C_2$ 则分别读出 $000_2=0$、$010_2=2$ 和 $011_2=3$。

\begin{figure}[htbp]
\centering
\begin{tabular}{c}
$0+0$\\[3pt]
\begin{quantikz}
\lstick{$\ket{0}$} & \qw      & \gate[6]{\text{两位加法器}} & \qw \\
\lstick{$\ket{0}$} & \qw      &                              & \meter{} \\
\lstick{$\ket{0}$} & \qw      &                              & \qw \\
\lstick{$\ket{0}$} & \qw      &                              & \qw \\
\lstick{$\ket{0}$} & \qw      &                              & \meter{} \\
\lstick{$\ket{0}$} & \qw      &                              & \meter{}
\end{quantikz}\\[14pt]
$1+1=01+01$\\[3pt]
\begin{quantikz}
\lstick{$\ket{0}$} & \gate{X} & \gate[6]{\text{两位加法器}} & \qw \\
\lstick{$\ket{0}$} & \gate{X} &                              & \meter{} \\
\lstick{$\ket{0}$} & \qw      &                              & \qw \\
\lstick{$\ket{0}$} & \qw      &                              & \qw \\
\lstick{$\ket{0}$} & \qw      &                              & \meter{} \\
\lstick{$\ket{0}$} & \qw      &                              & \meter{}
\end{quantikz}\\[14pt]
$1+2=01+10$\\[3pt]
\begin{quantikz}
\lstick{$\ket{0}$} & \gate{X} & \gate[6]{\text{两位加法器}} & \qw \\
\lstick{$\ket{0}$} & \qw      &                              & \meter{} \\
\lstick{$\ket{0}$} & \qw      &                              & \qw \\
\lstick{$\ket{0}$} & \qw      &                              & \qw \\
\lstick{$\ket{0}$} & \gate{X} &                              & \meter{} \\
\lstick{$\ket{0}$} & \qw      &                              & \meter{}
\end{quantikz}
\end{tabular}
\caption{\cref{fig:lec01-adder2} 的加法器针对三个输入进行配置（方框隐藏了其内部
门）：前导 $X$ 门加载操作数，三个仪表读取 $S_0,S_1,C_2$。自上而下各线为
$A_0,B_0,\ket{0},A_1,B_1,\ket{0}$。改编自 Brody，图 7.13、7.15 和 7.17。}
\label{fig:lec01-adder-runs}
\end{figure}

\begin{remark}[可逆性、辅助比特与垃圾位]
上述每个门都是基态的一个置换，所以加法器是一个大的置换酉算符 (unitary)——可逆的，
正如 \cref{thm:lec01-unitary-preserves} 对任意门所要求的那样。可逆性是用簿记换来的：
结果用 $\oplus$ 写到新的 $\ket{0}$ 辅助比特上，而不是覆写输入，所以线路同时也携带着
原始输入和一些中间``垃圾位''。Bennett 的技巧——用 $\CNOT$ 把最终答案复制出来，然后
把线路反向运行——会逆计算掉垃圾位并把辅助比特重置为 $\ket{0}$，这正是使加法器能够作为
一个干净子程序嵌入更大的相干算法中的原因。最后，众多的 Toffoli 门使加法器对真实硬件
颇为苛刻，在那里门误差随线路规模累积；这是一个早期的提示，说明为何随着线路增大，控制
噪声会变得至关重要。
\end{remark}
\section{态制备与测量}
\label{sec:lec01-prepmeas}

每一个量子算法都被两个并非门的操作所框定：它必须以\emph{制备} (preparing) 一个输入态开始，并以\emph{测量} (measuring) 一个输出态结束。\cref{sec:lec01-3} 的门模型描述了中间的酉演化；本节系统地处理这两个端点，建立在 \cref{sec:lec01-1} 已给出的计算基 Born 规则之上。两个端点都有代价，而这代价往往占据主导：一个漂亮的 $\poly(n)$ 门线路，如果加载其输入需要 $\bigo(2^{n})$ 个门，那它就毫无用处；而一个期望值也只不过与人们所能负担的测量重复次数一样好。

\subsection{从 \texorpdfstring{$\ket{0\cdots0}$}{|0...0>} 制备态}

按照约定，每个寄存器 (register) 都从全零态 $\ket{0}^{\otimes n}=\ket{0\cdots0}$ 开始。制备一个目标态意味着找到一个将 $\ket{0\cdots0}$ 映射到它的线路。

\begin{proposition}[基态与均匀叠加]\label{prop:lec01-prepbasic}
设 $x=x_{1}\cdots x_{n}\in\{0,1\}^{n}$。
\begin{enumerate}[label=\textup{(\alph*)}]
  \item 计算基态 $\ket{x}$ 可由 $\ket{0\cdots0}$ 制备，方法是对每个满足 $x_{j}=1$ 的量子比特 $j$ 施加 $X$；即
        $\bigl(\bigotimes_{j}X^{x_{j}}\bigr)\ket{0\cdots0}=\ket{x}$，至多使用
        $n$ 个门。
  \item 对所有 $2^{n}$ 个比特串的均匀叠加可由每个量子比特一个 Hadamard 制备，
        \begin{equation}\label{eq:lec01-uniform}
          H^{\otimes n}\ket{0\cdots0}
            =\frac{1}{\sqrt{2^{n}}}\sum_{x\in\{0,1\}^{n}}\ket{x},
        \end{equation}
        恰好使用 $n$ 个门。
\end{enumerate}
\end{proposition}

\begin{proof}
(a) $X^{x_{j}}\ket{0}=\ket{x_{j}}$ 逐比特作用，而各单量子比特结果的张量积为
$\ket{x_{1}}\cdots\ket{x_{n}}=\ket{x}$。
(b) 在单个量子比特上 $H\ket{0}=\tfrac{1}{\sqrt2}(\ket{0}+\ket{1})$。取张量幂并展开 $n$ 个和的乘积，
\[
  H^{\otimes n}\ket{0\cdots0}
    =\bigotimes_{j=1}^{n}\tfrac{1}{\sqrt2}(\ket{0}+\ket{1})
    =\frac{1}{\sqrt{2^{n}}}\sum_{x\in\{0,1\}^{n}}\ket{x},
\]
因为从每个因子中选取 $\ket{0}$ 或 $\ket{1}$ 恰好将所有比特串 $x$ 枚举一次。
\end{proof}

均匀叠加 \eqref{eq:lec01-uniform} 是第~I~部分查询算法的主力输入：一层深度为一的 Hadamard 将寄存器置于所有 $2^{n}$ 个输入的等权组合中，谕示 (oracle) 随后可在其上``并行''作用。

\subsection{单量子比特态合成}

一个一般的单量子比特目标，至多相差一个全局相位，是 \cref{prop:lec01-bloch} 的 Bloch 态
$\ket{\psi(\theta,\varphi)}=\cos\tfrac{\theta}{2}\ket{0}+\ee^{\ii\varphi}\sin\tfrac{\theta}{2}\ket{1}$。它可由 $\ket{0}$ 经一次绕 $y$ 轴的旋转 (rotation) 后再绕 $z$ 轴旋转而得到，其中旋转门为
\begin{equation}\label{eq:lec01-rotations}
  R_{y}(\theta)=\ee^{-\ii\theta Y/2}
    =\begin{pmatrix}\cos\tfrac{\theta}{2}&-\sin\tfrac{\theta}{2}\\[2pt]
      \sin\tfrac{\theta}{2}&\phantom{-}\cos\tfrac{\theta}{2}\end{pmatrix},
  \qquad
  R_{z}(\varphi)=\ee^{-\ii\varphi Z/2}
    =\begin{pmatrix}\ee^{-\ii\varphi/2}&0\\ 0&\ee^{\ii\varphi/2}\end{pmatrix}.
\end{equation}

\begin{proposition}[单量子比特制备]\label{prop:lec01-1qprep}
对任意 $\theta\in[0,\pi]$ 和 $\varphi\in[0,2\pi)$，
\[
  R_{z}(\varphi)\,R_{y}(\theta)\,\ket{0}
    =\ee^{-\ii\varphi/2}\Bigl(\cos\tfrac{\theta}{2}\ket{0}+\ee^{\ii\varphi}\sin\tfrac{\theta}{2}\ket{1}\Bigr),
\]
因此至多相差全局相位 $\ee^{-\ii\varphi/2}$，该线路用两个门制备出
$\ket{\psi(\theta,\varphi)}$。
\end{proposition}

\begin{proof}
由 \eqref{eq:lec01-rotations}，
$R_{y}(\theta)\ket{0}=\cos\tfrac{\theta}{2}\ket{0}+\sin\tfrac{\theta}{2}\ket{1}$。
施加对角的 $R_{z}(\varphi)$ 将 $\ket{0}$ 振幅乘以
$\ee^{-\ii\varphi/2}$，将 $\ket{1}$ 振幅乘以 $\ee^{\ii\varphi/2}$：
\[
  R_{z}(\varphi)R_{y}(\theta)\ket{0}
    =\ee^{-\ii\varphi/2}\cos\tfrac{\theta}{2}\ket{0}
     +\ee^{\ii\varphi/2}\sin\tfrac{\theta}{2}\ket{1}
    =\ee^{-\ii\varphi/2}\Bigl(\cos\tfrac{\theta}{2}\ket{0}+\ee^{\ii\varphi}\sin\tfrac{\theta}{2}\ket{1}\Bigr).
\]
整体因子 $\ee^{-\ii\varphi/2}$ 是一个不可观测的全局相位
（\cref{rem:lec01-global}）。
\end{proof}

\begin{example}[制备 $\ket{+\ii}$]\label{exa:lec01-prepplusi}
态 $\ket{+\ii}=\tfrac{1}{\sqrt2}(\ket{0}+\ii\ket{1})$ 的 Bloch 角为
$\theta=\pi/2$（等模）与 $\varphi=\pi/2$（相对相位 $\ii=\ee^{\ii\pi/2}$）。
由 \cref{prop:lec01-1qprep}，它可由 $R_{z}(\pi/2)R_{y}(\pi/2)\ket{0}$ 制备：
$R_{y}(\pi/2)$ 产生 $\ket{+}=\tfrac{1}{\sqrt2}(\ket{0}+\ket{1})$，而
$R_{z}(\pi/2)$ 随后将 $\ket{1}$ 振幅相对于 $\ket{0}$ 振幅乘以 $\ee^{\ii\pi/2}=\ii$，
至多相差全局相位地给出 $\tfrac{1}{\sqrt2}(\ket{0}+\ii\ket{1})$。注意同样有
$S\ket{+}=\ket{+\ii}$，所以这里单个 $S$ 门即可达到同样的结果。
\end{example}

\subsection{一般振幅编码}

从一个由 $2^{n}$ 个振幅 $\{a_{x}\}$ 构成的经典列表制备任意的 $n$ 量子比特态 $\ket{\psi}=\sum_{x\in\{0,1\}^{n}}a_{x}\ket{x}$，远比上述特殊情形昂贵得多。

\begin{theorem}[精确振幅编码的代价]\label{thm:lec01-prepcost}
每一个 $n$ 量子比特纯态 $\ket{\psi}=\sum_{x}a_{x}\ket{x}$ 都可由 $\ket{0\cdots0}$ 用单量子比特旋转和 $\CNOT$ 门构成的线路精确制备，且 $\bigo(2^{n})$ 个门就足够。反之，一个计数论证表明对一般态而言 $\Omega(2^{n})$ 个门是必要的，所以这个标度至多相差常数因子是最优的。
\end{theorem}

我们仅在这一层面陈述该结果，而不复现其构造；Möttönen、Vartiainen、Bergholm 与 Salomaa 的标准方法，以及等价的 Shende--Bullock--Markov 分解，通过一串均匀受控的 $R_{y}$ 与 $R_{z}$ 旋转来制备 $\ket{\psi}$，每个量子比特一层，将振幅的模与相位编码在旋转角之中。上界与一个信息论下界相匹配：指定 $\C^{2^{n}}$ 中一个一般的单位向量需要 $\sim 2^{n}$ 个实参数，而每个门只提供 $\bigo(1)$ 个，所以没有次指数线路能到达所有态。

\begin{remark}[高效的特殊情形]\label{rem:lec01-prepefficient}
\cref{thm:lec01-prepcost} 的指数代价是对\emph{所有}态取的最坏情形；实践中出现的态通常是有结构的，而结构换来效率。
\begin{itemize}[nosep]
  \item \emph{乘积态} $\ket{\psi_{1}}\otimes\cdots\otimes\ket{\psi_{n}}$ 耗费
        $\bigo(n)$ 个门：由 \cref{prop:lec01-1qprep} 独立地制备每个量子比特。
  \item \emph{可高效计算的振幅。} 若 $a_{x}$ 由某个可逆求值的公式给出
        （例如一个离散化的光滑密度，通过在条件概率上递归旋转的 Grover--Rudolph 方法），该态可用 $\poly(n)$ 个门制备。
  \item \emph{稀疏态}只有 $s$ 个非零振幅时，可用
        $\bigo(s\,n)$ 个门制备，与全维度 $2^{n}$ 无关。
\end{itemize}
当振幅来自外部数据而非公式时，人们转而假定一个黑箱式的\emph{态制备谕示}；这一输入模型及其在量子数值线性代数中的作用，将在 \cref{ch:lec06} 中展开。
\end{remark}

\subsection{投影测量}

\cref{sec:lec01-1} 的计算基 Born 规则是一个更一般测量公设的特例，我们现在完整地陈述它。

\begin{definition}[投影 (von Neumann) 测量]\label{def:lec01-projmeas}
一个\emph{可观测量} (observable) 是 $\Hilb$ 上的一个厄米算符 $O=O\dg$。将其谱分解写为 $O=\sum_{m}m\,P_{m}$，其中相异的实本征值 $m$ 标记相互正交的投影算符 $P_{m}$（满足 $P_{m}P_{m'}=\delta_{mm'}P_{m}$，$P_{m}\dg=P_{m}$）并分解恒等算符，$\sum_{m}P_{m}=\id$，那么对态 $\ket{\psi}$ 的相应\emph{投影测量} (projective measurement) 以如下概率返回结果 $m$
\begin{equation}\label{eq:lec01-projprob}
  p_{m}=\bra{\psi}P_{m}\ket{\psi}=\norm{P_{m}\ket{\psi}}^{2},
\end{equation}
并留下测量后的态
\begin{equation}\label{eq:lec01-projpost}
  \frac{P_{m}\ket{\psi}}{\sqrt{p_{m}}}\qquad(\text{当 }p_{m}>0\text{ 时有定义}).
\end{equation}
\end{definition}

这些概率非负且求和为一：由完备性，
$\sum_{m}p_{m}=\bra{\psi}\bigl(\sum_{m}P_{m}\bigr)\ket{\psi}=\braket{\psi|\psi}=1$。
选取在计算基中对角的 $O$，使得 $P_{m}=\proj{m}$ 是秩一的，就恰好恢复了
\cref{def:lec01-qubit} 的 Born 规则：$p_{m}=\bra{\psi}\proj{m}\ket{\psi}=\abs{\braket{m|\psi}}^{2}$，
且态坍缩到 $\ket{m}$。\Cref{def:lec01-projmeas} 在两方面推广了这一点——投影算符可以有秩 $>1$（一个\emph{简并}结果，如当我们只测量某些量子比特时），且本征基不必是计算基。

\begin{proposition}[期望值]\label{prop:lec01-expect}
可观测量 $O=\sum_{m}m\,P_{m}$ 在态 $\ket{\psi}$ 中的期望值为
\begin{equation}\label{eq:lec01-expect}
  \langle O\rangle \;=\; \sum_{m}m\,p_{m}\;=\;\sum_{m}m\,\bra{\psi}P_{m}\ket{\psi}\;=\;\bra{\psi}O\ket{\psi}.
\end{equation}
\end{proposition}

\begin{proof}
由结果随机变量均值的定义，$\langle O\rangle=\sum_{m}m\,p_{m}$。
代入 \eqref{eq:lec01-projprob} 并将求和移入内积之内，
$\sum_{m}m\,\bra{\psi}P_{m}\ket{\psi}=\bra{\psi}\bigl(\sum_{m}m\,P_{m}\bigr)\ket{\psi}=\bra{\psi}O\ket{\psi}$。
\end{proof}

\begin{example}[在 $\ket{0}$ 和 $\ket{+}$ 上测量 $Z$ 与 $X$]\label{exa:lec01-measZX}
取 $O=Z=(+1)\proj{0}+(-1)\proj{1}$，结果为 $\pm1$，投影算符为
$P_{+1}=\proj{0}$，$P_{-1}=\proj{1}$。在 $\ket{0}$ 上，结果以概率
$\abs{\braket{0|0}}^{2}=1$ 为 $+1$，所以 $\langle Z\rangle=1$。在
$\ket{+}=\tfrac{1}{\sqrt2}(\ket{0}+\ket{1})$ 上，每个结果的概率为 $\tfrac12$，
所以 $\langle Z\rangle=\tfrac12(+1)+\tfrac12(-1)=0$；等价地
$\bra{+}Z\ket{+}=0$。现在取 $O=X=(+1)\proj{+}+(-1)\proj{-}$。在 $\ket{0}$ 上，
$\abs{\braket{+|0}}^{2}=\abs{\braket{-|0}}^{2}=\tfrac12$，所以 $\langle X\rangle=0$；但
在 $\ket{+}$ 上，结果确定为 $+1$，且 $\langle X\rangle=1$。这两个可观测量具有互补的锐度：$\ket{+}$ 是 $X$ 的本征态（结果确定）但对 $Z$ 极度不确定，对 $\ket{0}$ 反之亦然。
\end{example}

\subsection{在任意基中测量}

真实硬件只在计算基中测量。要在某个其他标准正交基 $\{\ket{\phi_{m}}\}$ 中测量，首先用一个酉算符将该基旋转到计算基上，然后测量这些量子比特。

\begin{proposition}[测量前的基变换]\label{prop:lec01-basischange}
设 $\{\ket{\phi_{m}}\}_{m}$ 是一个标准正交基，并定义酉算符
$V=\sum_{m}\ket{m}\!\bra{\phi_{m}}$，它满足 $V\ket{\phi_{m}}=\ket{m}$。
在基 $\{\ket{\phi_{m}}\}$ 中测量 $\ket{\psi}$——即用投影算符
$\proj{\phi_{m}}$——产生结果 $m$ 的概率，与先施加
$V$ 再在计算基中测量 $V\ket{\psi}$ 的概率相同：
\[
  \abs{\braket{\phi_{m}|\psi}}^{2}=\abs{\braket{m|V|\psi}}^{2}.
\]
\end{proposition}

\begin{proof}
$V$ 是酉的，这由
$V\dg V=\sum_{m,m'}\ket{\phi_{m}}\braket{m|m'}\bra{\phi_{m'}}=\sum_{m}\proj{\phi_{m}}=\id$
得出，最后一个等式是该标准正交基的完备性
（\cref{thm:lec01-completeness}）。由定义
$\braket{m|V|\psi}=\sum_{m'}\braket{m|m'}\braket{\phi_{m'}|\psi}=\braket{\phi_{m}|\psi}$，
取模平方即得所述结论。
\end{proof}

\begin{example}[借助 $H$ 在 $X$ 基中测量]\label{exa:lec01-xbasis}
Hadamard 门满足 $H\ket{+}=\ket{0}$ 和 $H\ket{-}=\ket{1}$，所以
$V=H=\proj{0}H+\dots$ 恰是针对 $X$ 基 $\{\ket{+},\ket{-}\}$ 的基变换酉算符
$\ket{0}\!\bra{+}+\ket{1}\!\bra{-}$。因此对 $X$ 的测量可通过施加 $H$ 并在计算基中测量来实现，将 $\ket{+}\mapsto0$ 和 $\ket{-}\mapsto1$ 读出。例如，在
$\ket{\psi}=\cos\tfrac{\pi}{8}\ket{0}+\sin\tfrac{\pi}{8}\ket{1}$ 上，$X$ 基结果
$+$ 以概率
$\abs{\braket{+|\psi}}^{2}=\tfrac12\bigl(\cos\tfrac{\pi}{8}+\sin\tfrac{\pi}{8}\bigr)^{2}
=\cos^{2}\tfrac{\pi}{8}=\tfrac{2+\sqrt2}{4}\approx0.854$ 出现，所以
$\langle X\rangle=2\cos^{2}\tfrac{\pi}{8}-1=\cos\tfrac{\pi}{4}=\tfrac{1}{\sqrt2}$，与
$\bra{\psi}X\ket{\psi}=\sin\tfrac{\pi}{4}=\tfrac{1}{\sqrt2}$ 一致。
\end{example}

\subsection{广义测量与 POVM}

投影测量并非人们所能做的最一般之事：将系统耦合到一个辅助比特 (ancilla)，酉演化，然后对辅助比特作投影测量，就能在系统本身上实现一类更广的测量。这些由正算符值测度来刻画。

\begin{definition}[POVM]\label{def:lec01-povm}
一个\emph{正算符值测度} (positive operator-valued measure, POVM) 是 $\Hilb$ 上一组有限的算符
$\{E_{m}\}$，称为\emph{效应} (effects)，满足
\begin{equation}\label{eq:lec01-povm}
  E_{m}\succeq0\ \ \text{对所有 }m,\qquad \sum_{m}E_{m}=\id.
\end{equation}
作用于态 $\ket{\psi}$ 时，它以概率
$p_{m}=\bra{\psi}E_{m}\ket{\psi}$ 返回结果 $m$（更一般地，对密度算符 $\rho$ 为 $p_{m}=\Tr[E_{m}\rho]$）。
\end{definition}

每个 $E_{m}$ 的正性使 $p_{m}\ge0$，而完备性 $\sum_{m}E_{m}=\id$
使 $\sum_{m}p_{m}=1$，所以这些 $p_{m}$ 是一个真正的概率分布。投影测量是效应为正交投影算符的特例 $E_{m}=P_{m}$；一般的 POVM 舍弃了正交性与幂等性，允许比 $\Hilb$ 维数更多的结果，并且它本身不指定测量后的态。

\begin{theorem}[Naimark 扩张，陈述]\label{thm:lec01-naimark}
$\Hilb$ 上的每个 POVM $\{E_{m}\}$ 都源自更大空间上的一个投影测量：存在一个辅助空间、一个固定的辅助态、一个进入 $\Hilb\otimes\Hilb_{\mathrm{anc}}$ 的等距（或酉耦合），以及扩大空间上的正交投影算符
$\{\tilde P_{m}\}$，使得对每个 $\ket{\psi}$ 和每个 $m$，就扩张后的态
$\ket{\Psi}$ 而言有 $\bra{\psi}E_{m}\ket{\psi}=\bra{\Psi}\tilde P_{m}\ket{\Psi}$。
\end{theorem}

我们仅在陈述层面记录这一结果；一个构造性证明（将
$E_{m}=M_{m}\dg M_{m}$ 并由 $M_{m}$ 构建等距）在掌握量子信道之后是一个标准习题。它所证成的图景才是有用的：\emph{测量本身就是一个量子信道}，一个密度算符上的不可逆映射，而每个广义测量``其实''都是一个在辅助比特上被给予了舒展空间的投影测量。信道观点，以及量化任何 POVM 能将两个态区分得多好的距离度量，将在 \cref{ch:lec02} 中展开。

\subsection{从期望到估计}

测量产生一个随机结果，所以单次测量无法揭示一个期望值；人们重复制备-测量循环并求平均。若每次测量以 $\langle O\rangle$ 为均值返回 $\pm1$，那么 $M$ 次独立测量的经验平均的标准差为 $\bigo(1/\sqrt{M})$，所以将 $\langle O\rangle$ 估计到加性误差 $\varepsilon$ 需要 $M=\bigo(1/\varepsilon^{2})$ 次重复。这种对精度的平方依赖——\emph{散粒噪声} (shot-noise) 或\emph{标准量子极限}——是贯穿本课程的一个反复出现的代价；我们将其精确的统计处理（方差、Chernoff/Hoeffding 界，以及加速到 $\bigo(1/\varepsilon)$ 的平方振幅估计）推迟到 \cref{ch:lec02}，并将把 $\langle\psi|U|\psi\rangle$ 读作这样一个期望的线路——Hadamard 检验——推迟到 \cref{ch:lec04}。
\section{量子隐形传态与超密编码}
\label{sec:lec01-tele}

现在我们把前几节的各个部件——共享的贝尔对 (Bell pair)、$H$ 和 $\CNOT$ 门，以及测量——组装成两个协议，以展示纠缠 (entanglement) 究竟\emph{有何用}。隐形传态 (teleportation) 利用一个共享的贝尔对和两个经典比特传输一个未知量子比特；超密编码 (superdense coding) 则做相反的事，用一个量子比特和一个共享贝尔对发送两个经典比特。

\subsection{隐形传态协议}

Penelope 持有一个未知量子比特 $\ket{\psi}=a\ket{0}+b\ket{1}$，她希望把它传递给远方的 Odysseus。她可能并不知道 $a,b$，并且（由不可克隆定理 (no-cloning theorem)，\cref{sec:lec01-4}）无法复制该量子比特。他们共享的唯一资源是一个早先在他们还在一起时制备好的贝尔对 $\ket{\Phi^{+}}$：Penelope 保留其中一半，Odysseus 持有另一半。

把三个量子比特标记为 $1$（Penelope 的消息）、$2$（Penelope 持有的那一半对）以及 $3$（Odysseus 持有的那一半）。联合态为
\begin{equation}\label{eq:lec01-tele-init}
  \ket{\psi}_1\otimes\ket{\Phi^{+}}_{23}
    =\bigl(a\ket{0}+b\ket{1}\bigr)_1\otimes\tfrac{1}{\sqrt2}\bigl(\ket{00}+\ket{11}\bigr)_{23}.
\end{equation}

\begin{theorem}[隐形传态]\label{thm:lec01-tele}
Penelope 以量子比特~$1$（控制）对量子比特~$2$（目标）施加 $\CNOT$，再对量子比特~$1$ 施加 $H$，然后在计算基下测量量子比特~$1$ 和~$2$，得到两个比特 $(m_1,m_2)$，各自等概率出现。在给定 $(m_1,m_2)$ 的条件下，Odysseus 的量子比特~$3$ 为
\[
  X^{m_2}Z^{m_1}\ket{\psi}\quad\text{此外别无其他，}
\]
因此当 Odysseus 施加修正 $Z^{m_1}X^{m_2}$ 后，他的量子比特恰好为 $\ket{\psi}$。
\end{theorem}

\begin{proof}
将 \eqref{eq:lec01-tele-init} 对量子比特 $2,3$ 展开：
\[
  \tfrac{1}{\sqrt2}\bigl(a\ket{0}_1(\ket{00}+\ket{11})_{23}
      +b\ket{1}_1(\ket{00}+\ket{11})_{23}\bigr).
\]
施加 $\CNOT_{1\to2}$，它恰在量子比特~$1$ 为 $\ket{1}$ 时翻转量子比特~$2$：
\[
  \tfrac{1}{\sqrt2}\bigl(a\ket{0}_1(\ket{00}+\ket{11})_{23}
      +b\ket{1}_1(\ket{10}+\ket{01})_{23}\bigr).
\]
现对量子比特~$1$ 施加 $H$，利用 $H\ket{0}=\tfrac{1}{\sqrt2}(\ket{0}+\ket{1})$ 和
$H\ket{1}=\tfrac{1}{\sqrt2}(\ket{0}-\ket{1})$：
\[
  \tfrac12\Bigl(a(\ket{0}+\ket{1})_1(\ket{00}+\ket{11})_{23}
      +b(\ket{0}-\ket{1})_1(\ket{10}+\ket{01})_{23}\Bigr).
\]
按量子比特 $1,2$（Penelope 的寄存器）的四个取值重新分组，提出量子比特~$3$：
\begin{equation}\label{eq:lec01-tele-grouped}
\begin{aligned}
  \tfrac12\Bigl(\;
    &\ket{00}_{12}\,(a\ket{0}+b\ket{1})_3
    +\ket{01}_{12}\,(a\ket{1}+b\ket{0})_3\\
    +\;&\ket{10}_{12}\,(a\ket{0}-b\ket{1})_3
    +\ket{11}_{12}\,(a\ket{1}-b\ket{0})_3\;\Bigr).
\end{aligned}
\end{equation}
这四项具有相等的权重 $\tfrac12$，因此每个结果 $(m_1,m_2)$ 以概率 $\tfrac14$ 出现。读出每个分支中的量子比特~$3$：
\[
  \begin{array}{llll}
    (0,0):&\ket{\psi}, & (0,1):&X\ket{\psi},\\
    (1,0):&Z\ket{\psi}, & (1,1):&XZ\ket{\psi} \ (\text{至多相差一个全局相位}),
  \end{array}
\]
其中我们用到 $X\ket{\psi}=a\ket{1}+b\ket{0}$ 及 $Z\ket{\psi}=a\ket{0}-b\ket{1}$。
在每个分支中 Odysseus 的态都是 $X^{m_2}Z^{m_1}\ket{\psi}$；由于 $X,Z$ 都是自身的逆，施加 $Z^{m_1}X^{m_2}$（即若 $m_2=1$ 则施加 $X$，再若 $m_1=1$ 则施加 $Z$）即可把量子比特~$3$ 还原为 $\ket{\psi}$。
\end{proof}

\begin{remark}[为何隐形传态需要经典通信，却不能超越它]
协议中有四个特征值得细究。
\begin{itemize}[nosep]
  \item \emph{经典比特不可或缺。} 在 Odysseus 获知 $(m_1,m_2)$ 之前，他的量子比特在对 Penelope 的随机结果取平均的意义下是最大混合态 (maximally mixed state) $\tfrac12\id$（$\ket{\psi}$ 的四个修正像的平均），不携带任何关于 $\ket{\psi}$ 的信息。只有通过普通信道发送的那两个经典比特，才能告诉 Odysseus 该施加哪个修正。
  \item \emph{不存在超光速信号。} 由于无论 Penelope 做什么，Odysseus 的边缘态都是 $\tfrac12\id$，所以 Penelope 对量子比特~$1,2$ 的任何操作都不会改变 Odysseus 在本地所能观测到的统计。只有当经典消息到达时，态才被完成，而该消息的传播不快于光速。这就是\emph{不可超光速通信} (no-signaling) 原理，它正是 \cref{sec:lec01-chsh} 中的非定域 (nonlocal) 关联能与相对论和平共存的原因。
  \item \emph{源被摧毁。} Penelope 的测量把量子比特~$1$ 坍缩到一个计算基态，从原件中抹去了 $a,b$。隐形传态\emph{移动}态而非复制它——这与不可克隆定理（\cref{thm:lec01-nocloning}）一致，否则一个端到端的映射 $\ket{\psi}\mapsto\ket{\psi}\otimes\ket{\psi}$ 将违背它。
  \item \emph{传递的是信息，而非物质。} 跨越距离的是量子\emph{信息}——振幅 (amplitude) $a,b$——它们从两个经典比特连同共享的纠缠在 Odysseus \emph{预先存在}的量子比特上被重建；没有任何物理量子比特或质量被传输。（Brody，\emph{The Joy of Quantum Computing}，第~4 章，把这讲成 Penelope 利用他们分别前创建的纠缠对，将一个量子比特隔着爱琴海隐形传态给远航的 Odysseus。）
\end{itemize}
\end{remark}

\subsection{超密编码}

把同一个贝尔资源``反向''运行，使 Penelope 能够通过物理传输\emph{单个}量子比特来发送\emph{两个}经典比特。

\begin{proposition}[超密编码]\label{prop:lec01-superdense}
Penelope 和 Odysseus 共享 $\ket{\Phi^{+}}=\tfrac{1}{\sqrt2}(\ket{00}+\ket{11})$（量子比特~$1$ 属 Penelope，量子比特~$2$ 属 Odysseus）。为发送两个比特 $(x,y)\in\{0,1\}^2$，Penelope 对她的量子比特施加 $Z^{x}X^{y}$ 并发送给 Odysseus。Odysseus 随后施加 $\CNOT_{1\to2}$，再对量子比特~$1$ 施加 $H$ 并测量两个量子比特，以确定性地恢复 $(x,y)$。
\end{proposition}

\begin{proof}
这四个局域操作把 $\ket{\Phi^{+}}$ 映射到四个标准正交的贝尔态：
\[
  \begin{aligned}
    (0,0)\colon\ \id\ket{\Phi^{+}}&=\tfrac{1}{\sqrt2}(\ket{00}+\ket{11}),
    & (0,1)\colon\ X_1\ket{\Phi^{+}}&=\tfrac{1}{\sqrt2}(\ket{10}+\ket{01}),\\
    (1,0)\colon\ Z_1\ket{\Phi^{+}}&=\tfrac{1}{\sqrt2}(\ket{00}-\ket{11}),
    & (1,1)\colon\ Z_1X_1\ket{\Phi^{+}}&=\tfrac{1}{\sqrt2}(\ket{10}-\ket{01}).
  \end{aligned}
\]
这四个态相互正交，因此可以完美区分。映射 $H_1\,\CNOT_{1\to2}$ 恰好是 \cref{thm:lec01-bellprep} 中贝尔制备线路的逆，所以它分别把这四个态送到 $\ket{00},\ket{01},\ket{10},\ket{11}$（可如该证明中那样逐一验证每种情形）。因此在计算基下测量确定性地返回 $(x,y)$。
\end{proof}

\begin{remark}[资源核算]
隐形传态与超密编码是对偶的：一个量子比特的通信加上一个共享的纠缠对，价值相当于两个经典比特，反之亦然。两个协议都不能让纠缠单独传输信息——每个都需要一个实际信道（两个经典比特，或一个量子比特）——这再次是不可超光速通信原理在起作用。
\end{remark}

\section{不可克隆定理与可逆性}
\label{sec:lec01-4}

在经典情形下，复制一个比特是平凡的。而在量子力学中，复制一个\emph{未知}态是不可能的，这一事实由 Wootters、Zurek 和 Dieks 于 1982 年发现。一个复制器将是这样一个门 $U$：给定待复制的态和一个固定的空白寄存器 $\ket{s}$，输出两份副本。

\begin{theorem}[不可克隆定理]\label{thm:lec01-nocloning}
不存在 $\Hilb\otimes\Hilb$ 上的酉算符 (unitary operator) $U$ 和空白态 $\ket{s}\in\Hilb$ 使得
\[
  U\big(\ket{\psi}\otimes\ket{s}\big)=\ket{\psi}\otimes\ket{\psi}
  \qquad\text{对每个归一化的 } \ket{\psi}\in\Hilb,
\]
除非 $\dim\Hilb=1$。更确切地说，任何这样的 $U$ 能同时克隆两个态 $\ket{x_{1}},\ket{x_{2}}$，仅当它们相等或正交。
\end{theorem}

\begin{proof}
假设这样的 $U$ 和 $\ket{s}$ 存在。将克隆方程应用于两个归一化态 $\ket{x_{1}}$ 和 $\ket{x_{2}}$：
\[
  U\ket{x_{1}}\ket{s}=\ket{x_{1}}\ket{x_{1}},\qquad
  U\ket{x_{2}}\ket{s}=\ket{x_{2}}\ket{x_{2}}.
\]
取这两个方程的内积 (inner product)。左边，由于 $U$ 是酉的，它保持内积（\cref{thm:lec01-unitary-preserves}），且 $\braket{s|s}=1$，故
\[
  \big(\bra{x_{1}}\bra{s}\big)U\dg U\big(\ket{x_{2}}\ket{s}\big)
   = \braket{x_{1}|x_{2}}\,\braket{s|s}
   = \braket{x_{1}|x_{2}}.
\]
右边，$\ket{x_{1}}\ket{x_{1}}$ 与 $\ket{x_{2}}\ket{x_{2}}$ 的内积在张量积 (tensor product) 上分解，
\[
  \big(\bra{x_{1}}\bra{x_{1}}\big)\big(\ket{x_{2}}\ket{x_{2}}\big)
   = \braket{x_{1}|x_{2}}\,\braket{x_{1}|x_{2}}
   = \braket{x_{1}|x_{2}}^{2}.
\]
令两边相等，$\braket{x_{1}|x_{2}} = \braket{x_{1}|x_{2}}^{2}$。记 $z=\braket{x_{1}|x_{2}}$，此式即为 $z=z^{2}$，亦即 $z(z-1)=0$，故 $z\in\{0,1\}$。于是 $\ket{x_{1}}$ 和 $\ket{x_{2}}$ 要么相同（$z=1$），要么正交（$z=0$）。一个通用克隆器需要对\emph{所有}对都成立，但维数 $\geq 2$ 的空间包含内积严格介于 $0$ 和 $1$ 之间的态（例如 $\ket{0}$ 和 $\ket{+}$，其 $\braket{0|+}=\tfrac{1}{\sqrt2}$）。这一矛盾证明了不存在通用克隆器；只有相互正交的态才能被克隆。
\end{proof}

证明的本质在于：克隆 $\ket{\psi}\mapsto\ket{\psi}\otimes\ket{\psi}$ 是 $\ket{\psi}$ 的一个非线性（二次）映射，而每个量子门都是线性的；线性性迫使内积约束 $z=z^{2}$，这是任何连续统的态都无法满足的。

\begin{remark}[若干推论]\leavevmode
\begin{itemize}[nosep]
  \item \emph{不可广播。} 既然未知态不能被复制，它也就不能被广播给多方。
  \item \emph{$\CNOT$ 不克隆叠加态。} 对于 $\ket{x}=a\ket{0}+b\ket{1}$，
        \[
          \CNOT\ket{x}\ket{0}=a\ket{00}+b\ket{11}\neq\ket{x}\ket{x}
          =\ket{x}\otimes(a\ket{0}+b\ket{1}),
        \]
        因为右边包含交叉项 $ab\,\ket{01},\,ab\,\ket{10}$，而左边没有。$\CNOT$ 只复制计算基的信息。
  \item \emph{BB84 安全性。} 由于窃听者既不能复制传输的量子比特，又不能在不扰动非正交态的情况下测量它们，BB84 协议中的截获必然引入可检测的错误，这赋予量子密钥分发以安全性。
\end{itemize}
\end{remark}

\begin{remark}[不可克隆保护了不可超光速通信]\label{rem:lec01-clonesignal}
克隆不仅仅是不便——若它可行，就会让纠缠以超光速发送信号，这与相对论直接矛盾。假设 Penelope 和 Odysseus 共享 $\ket{\Phi^{+}}$，而 Penelope 把一个比特编码进她的\emph{测量基选择}：要发送 $0$ 她在计算（$z$）基下测量，要发送 $1$ 则在 Hadamard（$x$）基下测量。她的测量相应地把 Odysseus 的量子比特留在 $\{\ket{0},\ket{1}\}$ 或 $\{\ket{+},\ket{-}\}$ 中，但如 \cref{ex:lec01-bellreduced} 所示，Odysseus 的约化态 (reduced state) 无论哪种情形都是 $\tfrac12\id$，所以单次测量无法揭示 Penelope 的基。\emph{假如} Odysseus 能把他的量子比特克隆成许多份，他就能在统计上估计其态：一个在 $z$ 和 $x$ \emph{两种}基下都是 $50/50$ 的系综（即 $\tfrac12\id$ 情形），与一个在 Penelope 所用的那种基下尖锐集中的系综是可区分的，这会在任意距离上瞬时揭示 Penelope 的比特。不可克隆定理恰恰禁止了这种态估计，从而封堵了这条本应存在的超光速信道——这与 \cref{sec:lec01-tele} 中隐形传态所建设性地得出的不可超光速通信结论相同。
\end{remark}

同样的内积论证，反向运行，禁止了对未知态的删除。

\begin{theorem}[不可删除定理]\label{thm:lec01-nodeleting}
在维数 $\geq 2$ 的希尔伯特空间上，不存在酉算符 $U$ 和固定目标 $\ket{\Sigma}$ 使得
\[
  U\big(\ket{\psi}\otimes\ket{\psi}\big)=\ket{\psi}\otimes\ket{\Sigma}
  \qquad\text{对所有归一化的 } \ket{\psi}.
\]
\end{theorem}

\begin{proof}
假设这样的 $U$ 存在，并把它应用于两个态 $\ket{x_{1}},\ket{x_{2}}$。取内积，左边用酉性、右边作分解，
\[
  \braket{x_{1}|x_{2}}^{2}=\braket{x_{1}|x_{2}}\,\braket{\Sigma|\Sigma}
                          =\braket{x_{1}|x_{2}},
\]
因为 $\braket{\Sigma|\Sigma}=1$。同样，$z=\braket{x_{1}|x_{2}}$ 满足 $z^{2}=z$，故 $z\in\{0,1\}$，这对非正交态不成立。因此不存在通用的删除操作。
\end{proof}

\paragraph{可逆性。} 由于每个门 $U$ 都是酉的，它有一个逆 $U^{-1}=U\dg$，其本身也是一个门：任何由序列 $U_m\cdots U_1$ 实现的计算，都可以通过施加 $U_1\dg\cdots U_m\dg$ 来撤销。因此量子演化是\emph{逻辑可逆的}，这与经典的不可逆门（如 \textsc{and}，它丢弃一个比特并按 Landauer 原理耗散能量）形成鲜明对比。正是这种可逆性，使得经典子程序必须嵌入上一节的 XOR-谕示 (oracle) 形式 $\ket{x}\ket{b}\mapsto\ket{x}\ket{b\oplus f(x)}$ 中，也使得\emph{反计算} (uncomputation)——将产生辅助比特 (ancilla) 的门反向运行以把草稿寄存器重置为 $\ket{0}$——成为量子算法设计中的常规一环。唯一本质上不可逆的步骤是测量，它坍缩态并摧毁测量前的叠加。
\section{密度算符与混合态}
\label{sec:lec01-5}

迄今为止，每个态都是单个单位向量 $\ket{\psi}$，称为\emph{纯态 (pure state)}。然而我们常常只知道系统以经典概率 $p_i$ 处于若干个态 $\ket{\psi_i}$ 之一（一个\emph{系综 (ensemble)} $\{p_i,\ket{\psi_i}\}$），或者我们只持有一个更大纠缠系统的一部分。这两种情形都可以由密度算符 (density operator) 形式体系来刻画。

\begin{definition}[密度算符]\label{def:lec01-density}
系综 $\{p_i,\ket{\psi_i}\}$（其中 $p_i\ge0$，$\sum_i p_i=1$）的\emph{密度算符}是
\begin{equation}\label{eq:lec01-density}
  \rho\;=\;\sum_i p_i\,\proj{\psi_i}.
\end{equation}
纯态 $\ket{\psi}$ 对应于秩一的 $\rho=\proj{\psi}$。
\end{definition}

对任意带有投影算符 $\{P_m\}$ 的可观测量测量，其结果概率以及 $M=\sum_m m\,P_m$ 的期望 (expectation) 为
\begin{equation}\label{eq:lec01-bornrho}
  p(m)=\Tr[P_m\rho],\qquad \E[M]=\Tr[M\rho],
\end{equation}
对纯态而言，由于 $\Tr[P_m\proj{\psi}]=\bra{\psi}P_m\ket{\psi}$，这归结为 \cref{sec:lec01-1} 的玻恩规则。

\begin{proposition}[密度算符的刻画]\label{prop:lec01-densitychar}
算符 $\rho$ 是某个系综的密度算符，当且仅当
\begin{enumerate}[label=(\roman*)]
  \item $\Tr\rho=1$，并且
  \item $\rho\succeq0$（厄米且半正定）。
\end{enumerate}
\end{proposition}

\begin{proof}
（$\Rightarrow$）由 \eqref{eq:lec01-density}，$\rho\dg=\sum_i p_i(\proj{\psi_i})\dg=\rho$
是厄米的；对任意 $\ket{\phi}$，$\bra{\phi}\rho\ket{\phi}=\sum_i p_i\abs{\braket{\phi|\psi_i}}^{2}\ge0$，
故 $\rho\succeq0$；并且 $\Tr\rho=\sum_i p_i\Tr[\proj{\psi_i}]=\sum_i p_i=1$。
（$\Leftarrow$）一个半正定的厄米算符具有谱分解
$\rho=\sum_i\lambda_i\proj{\lambda_i}$，其中 $\lambda_i\ge0$，且由
$\Tr\rho=1$ 有 $\sum_i\lambda_i=1$；这把 $\rho$ 表示为系综
$\{\lambda_i,\ket{\lambda_i}\}$ 的密度算符。
\end{proof}

\begin{definition}[纯度；纯态与混合态]\label{def:lec01-purity}
$\rho$ 的\emph{纯度 (purity)} 是 $\Tr[\rho^{2}]\in(0,1]$。一个态是\emph{纯的}当且仅当
$\rho^{2}=\rho$（等价地 $\Tr[\rho^{2}]=1$，等价地 $\rho$ 是秩一投影算符）；
否则 $\Tr[\rho^{2}]<1$，且 $\rho$ 称为\emph{混合态 (mixed state)}。
\end{definition}

例如，最大混合的单量子比特态 $\rho=\tfrac12\id$ 具有
$\Tr[\rho^{2}]=\tfrac12$：它描述一个以 $\tfrac12$ 的概率各为 $\ket{0}$ 或 $\ket{1}$ 的量子比特，
并且没有任何测量能把它与 $\ket{+}$ 和 $\ket{-}$ 的均匀混合区分开——系综分解并不唯一，只有
$\rho$ 是物理的。

\subsection{偏迹与约化态}

为了描述复合系统的一部分，我们把其余部分“迹掉”。

\begin{definition}[偏迹]\label{def:lec01-ptrace}
设 $\Hilb=\Hilb_A\otimes\Hilb_B$，其中 $\{\ket{j}_B\}$ 是 $\Hilb_B$ 的标准正交基。
\emph{对 $B$ 的偏迹 (partial trace)} 是把 $\Hilb$ 上的算符
$M$ 映到 $\Hilb_A$ 上算符的线性映射
\begin{equation}
  \Tr_B[M]\;=\;\sum_j (\id_A\otimes\bra{j}_B)\,M\,(\id_A\otimes\ket{j}_B).
\end{equation}
在乘积算符上它的作用为 $\Tr_B[A\otimes B]=A\,\Tr[B]$。子系统 $A$ 的\emph{约化密度算符 (reduced
density operator)} 为 $\rho_A=\Tr_B[\rho]$；它对仅在 $A$ 上进行的每个测量都给出正确的统计
\eqref{eq:lec01-bornrho}。
\end{definition}

\begin{example}[贝尔对的约化态]\label{ex:lec01-bellreduced}
取贝尔态 $\ket{\Phi^{+}}=\tfrac{1}{\sqrt2}(\ket{00}+\ket{11})$，其密度算符为
\[
  \rho=\proj{\Phi^{+}}
   =\tfrac12\big(\ket{00}+\ket{11}\big)\big(\bra{00}+\bra{11}\big)
   =\tfrac12\big(\proj{00}+\ket{00}\!\bra{11}+\ket{11}\!\bra{00}+\proj{11}\big).
\]
迹掉第二个量子比特。利用
$\Tr_B[\ket{ab}\!\bra{cd}]=\ket{a}\!\bra{c}\,\braket{d|b}$（因此只有 $b=d$ 的项
存留），
\[
  \rho_A=\Tr_B[\rho]
   =\tfrac12\big(\proj{0}\,\underbrace{\braket{0|0}}_{1}+\ket{0}\!\bra{1}\,\underbrace{\braket{1|0}}_{0}
   +\ket{1}\!\bra{0}\,\underbrace{\braket{0|1}}_{0}+\proj{1}\,\underbrace{\braket{1|1}}_{1}\big)
   =\tfrac12\big(\proj{0}+\proj{1}\big)=\tfrac12\id.
\]
于是贝尔对一半的约化态是\emph{最大混合的}，其纯度为
$\Tr[\rho_A^{2}]=\tfrac12<1$，尽管全局态是纯的
（$\Tr[\rho^{2}]=1$）。缺失的信息完全存在于两个量子比特之间的
关联中：这种限制到子系统时纯度的损失是纠缠的
标志。
\end{example}

酉演化通过共轭作用于密度算符，
$\rho\mapsto U\rho U\dg$，对纯态而言它重现了
$\ket{\psi}\mapsto U\ket{\psi}$。然而并非每个物理操作都是酉的；
一般的（完全正、保迹的）量子信道——以及比较密度算符的
距离度量——是下一讲的主题。本课程的大部分内容我们都使用纯态和酉操作，
但只要讨论子系统、噪声或统计混合，密度算符仍是恰当的工具。

\section{量子算法的复杂度}
\label{sec:lec01-6}

最后我们确定评判量子算法所用的标尺。设 $n$ 为表示输入所需的量子比特数。

\paragraph{门复杂度。} 一条线路的\emph{门复杂度 (gate complexity)} 是它所包含的
初等（单比特和双比特）门的总数。一个量子算法是\emph{高效的}，如果它的门复杂度为
$\bigo(\poly(n))$。由于测量结果是概率性的，只要算法以足够高的概率 $p$ 给出正确答案，我们就满意：
对判定问题我们要求 $p>2/3$
（或者仅要求 $p>1/2+1/\poly(n)$），并把线路重复 $\bigo(\log(1/\delta))$
次再取多数表决，根据切尔诺夫界 (Chernoff bound) 可把误差压到任意 $\delta>0$ 之下。

\paragraph{查询复杂度。} 输入常常由一个黑箱酉算符——一个
\emph{量子谕示 (quantum oracle)}——提供，我们只计对它的调用次数，而忽略周围门的
代价。对布尔函数 $f:\{0,1\}^{n}\to\{0,1\}$，两个标准谕示是
异或（或比特）谕示 $\Uf\ket{x}\ket{b}=\ket{x}\ket{b\oplus f(x)}$
以及相位谕示
\begin{equation}
  \Uf\ket{x}=(-1)^{f(x)}\ket{x},
\end{equation}
后者以\emph{相位回踢}的方式返回答案。查询复杂度 (query complexity) 是
第~I~部分各算法（Deutsch--Jozsa、Grover、相位估计）的
核心代价度量：量子加速最简洁的表述就是对谕示调用次数的减少。注意，拥有对 $f$
的查询权限并不意味着我们在全局上知道 $f$——找到一个满足 $f(x)=0$ 的 $x$ 仍可能很困难。

\paragraph{线路深度。} 把门分组成由相互对易、作用于不相交量子比特的门构成的层
（\cref{def:lec01-depth}）；\emph{深度 (depth)} 是层数，即最长的输入到输出路径。由于同一层中的门并行运行，
深度是量子版本的墙上时钟时间。它也是近期硬件的
约束性物理限制，这些硬件的量子比特只在有限的\emph{相干时间 (coherence time)} 内保持相干；
人们常常用更大的门数或更多的
重复，去换取一条能装进相干预算的更浅线路。

\paragraph{空间与辅助比特。} 一个算法的量子比特分为\emph{系统}
寄存器（保存关注的数据）和\emph{辅助 (ancilla)} 寄存器（为使操作酉化所需的
草稿空间）。只要可能，人们就把辅助比特\emph{反计算 (uncompute)}
回到 $\ket{0}$ 以便复用；可以这样释放的寄存器是
\emph{工作寄存器}，在渐近分析中它的代价常常不被计入。

\paragraph{渐近记号。} 我们对上述所有代价都按 $n$（以及诸如 $1/\varepsilon$ 的精度参数）
渐近地度量。回忆 $g(n)=\bigo(h(n))$
意味着对某个常数 $C>0$ 和所有大的 $n$ 有 $\abs{g(n)}\le C\,h(n)$；我们用
$\poly(n)$ 表示 $\bigo(n^{c})$，用 $\polylog(n)$ 表示 $\poly(\log n)$。
一个\emph{指数}加速把经典的 $\bigo(2^{n})$ 代价换成量子的
$\poly(n)$ 代价，而 Grover 算法给出\emph{平方}加速，
$\bigo(2^{n})\to\bigo(2^{n/2})$。这些就是我们将贯穿始终进行的比较。

\exercisesection
\begin{exercise}\label{ex:lec01-1}
设 $\ket{\psi}=2\ii\ket{0}-\ket{1}\in\C^{2}$（未归一化）。求 $\norm{\psi}$ 和
归一化态 $\ket{\hat\psi}$，并给出在计算基中测得 $0$
和 $1$ 的概率。
\end{exercise}

\begin{exercise}\label{ex:lec01-2}
证明 $\ket{+}=\tfrac{1}{\sqrt2}(\ket{0}+\ket{1})$ 和
$\ket{-}=\tfrac{1}{\sqrt2}(\ket{0}-\ket{1})$ 构成 $\C^{2}$ 的一组标准正交基，
并把 $\ket{0}$ 和 $\ket{1}$ 用这组基表示。
\end{exercise}

\begin{exercise}\label{ex:lec01-3}
利用柯西--施瓦茨不等式的等号情形
（\cref{thm:lec01-cs}），证明两个\emph{单位}向量 $\ket{u},\ket{v}$ 满足
$\abs{\braket{u|v}}=1$ 当且仅当对某个实数 $\gamma$ 有 $\ket{v}=\ee^{\ii\gamma}\ket{u}$——
即它们是同一个物理态。
\end{exercise}

\begin{exercise}\label{ex:lec01-4}
设 $\ket{u}$ 为单位向量，$P=\proj{u}$。证明 $P^{2}=P$ 和 $P\dg=P$，
并计算 $\Tr(P)$。
\end{exercise}

\begin{exercise}\label{ex:lec01-5}
求态
$\ket{\psi}=\tfrac{1}{2}\ket{0}+\tfrac{\sqrt3}{2}\ee^{\ii\pi/3}\ket{1}$ 的布洛赫角 $(\theta,\varphi)$。
\end{exercise}

\begin{exercise}\label{ex:lec01-6}
判断 $\ket{\psi}=\tfrac{1}{\sqrt2}(\ket{00}+\ket{01})$ 是否为乘积态。
若是，给出一个明确的因式分解；若否，证明它是纠缠的。
\end{exercise}

\begin{exercise}\label{ex:lec01-7}
通过矩阵乘法验证 $H$ 是酉的且自逆
（$H\dg=H$ 且 $H^{2}=I$）。然后利用 \cref{lem:lec01-conj} 和
\cref{prop:lec01-yxz} 计算 $HYH$。
\end{exercise}

\begin{exercise}\label{ex:lec01-8}
计算 $\CNOT\ket{+}\ket{+}$ 并判断结果是否纠缠。然后
给出一个酉算符 $U$ 和空白态 $\ket{s}$，使得 $U\ket{0}\ket{s}=\ket{0}\ket{0}$
且 $U\ket{1}\ket{s}=\ket{1}\ket{1}$，并解释为什么这不违反
不可克隆定理（\cref{thm:lec01-nocloning}）。
\end{exercise}

\begin{exercise}\label{ex:lec01-9}
考虑双量子比特态 $\ket{\psi}=\tfrac{1}{\sqrt2}(\ket{00}+\ket{11})$，但现在
只测量\emph{第一个}量子比特。利用密度算符，(a) 计算
约化态 $\rho_A=\Tr_B[\proj{\psi}]$ 及其纯度 $\Tr[\rho_A^{2}]$；
(b) 给出第一个量子比特得到结果 $0$ 的概率，以及由此得到的第二个量子比特的
（归一化）态。与乘积态 $\ket{+}\ket{+}$ 的约化态作比较。
\end{exercise}

\begin{exercise}\label{ex:lec01-10}
在隐形传态协议中（\cref{thm:lec01-tele}），假设 Penelope 的测量
得到 $(m_1,m_2)=(1,0)$。(a) 写出 Odysseus 校正前的态并验证它
等于 $Z\ket{\psi}$。(b) 说出 Odysseus 施加的校正并确认它恢复了
$\ket{\psi}$。(c) 在经典比特到达\emph{之前}对全部四个等可能结果求平均，
证明 Odysseus 的约化态是 $\tfrac12\id$，并用
一句话解释为什么这正是禁止超光速通信的原因。
\end{exercise}

\begin{exercise}\label{ex:lec01-11}
验证超密编码的一个情形（\cref{prop:lec01-superdense}）。从
$\ket{\Phi^{+}}$ 出发，对量子比特~$1$ 施加 $Z^{1}X^{1}=Z X$ 以编码 $(x,y)=(1,1)$，然后
对量子比特~$1$ 施加 $\CNOT_{1\to2}$ 再施加 $H$。证明一次计算基
测量以确定性返回 $11$。（固定全局相位约定 $ZX\ket{0}
=\ket{1}$，$ZX\ket{1}=-\ket{0}$。）
\end{exercise}

\begin{exercise}\label{ex:lec01-12}
利用 \cref{prop:lec01-chsh-corr} 的量子关联函数 $\E[AB]=\cos(\alpha-\beta)$，
考虑 CHSH 设置 $\alpha_1=0$，
$\alpha_2=\tfrac{\pi}{2}$，$\beta_1=\tfrac{\pi}{4}$，$\beta_2=\tfrac{3\pi}{4}$。
(a) 计算这四个关联函数并确认 $\E[S]=2\sqrt2$。(b) 通过最大化
$f(\alpha_1,\alpha_2,\beta_1,\beta_2)=\cos(\alpha_1-\beta_1)-\cos(\alpha_1-\beta_2)
+\cos(\alpha_2-\beta_1)+\cos(\alpha_2-\beta_2)$，证明对这个态而言没有任何角度选择
能使 $\abs{\E[S]}$ 超过 $2\sqrt2$。\emph{提示：}先固定
$\beta_1,\beta_2$ 并对 $\alpha_1,\alpha_2$ 优化，把两个
依赖 $\alpha$ 的对各写成形如 $R\cos(\alpha-\phi)$ 的单个余弦。
\end{exercise}

\begin{exercise}\label{ex:lec01-13}
给出明确的旋转角，使用 \cref{prop:lec01-1qprep}，以
$R_{z}(\varphi)R_{y}(\theta)\ket{0}$ 的形式从 $\ket{0}$ 制备态
$\ket{\psi}=\tfrac{1}{2}\ket{0}+\tfrac{\sqrt3}{2}\ee^{\ii\pi/2}\ket{1}$。
给出 $\theta$ 和 $\varphi$ 的值，并验证所得态在相差一个全局相位的意义下与
$\ket{\psi}$ 一致。
\end{exercise}

\begin{exercise}\label{ex:lec01-14}
设 $\ket{\psi}=\cos\tfrac{\pi}{8}\ket{0}+\sin\tfrac{\pi}{8}\ket{1}$，并设 $O$ 为
“沿角度 $\vartheta=\tfrac{\pi}{4}$”测量的可观测量，即以
$\ket{\vartheta}=\cos\tfrac{\vartheta}{2}\ket{0}+\sin\tfrac{\vartheta}{2}\ket{1}$
（结果 $+1$）及其正交伴（结果 $-1$）为本征态的厄米算符。(a) 写出两个投影算符
$P_{\pm1}$。(b) 计算 $\ket{\psi}$ 的结果概率 $p_{\pm1}$ 和
期望 $\langle O\rangle$。(c) 指出 \cref{prop:lec01-basischange} 中把它化为计算基测量的
基变换酉算符 $V$。
\end{exercise}

\begin{exercise}\label{ex:lec01-15}
考虑 $\C^{2}$ 上的三个算符
$E_{0}=\tfrac12\proj{0}$、$E_{1}=\tfrac12\proj{1}$ 和 $E_{2}=\tfrac12\id$。
(a) 验证 $\{E_{0},E_{1},E_{2}\}$ 是一个有效的 POVM（\cref{def:lec01-povm}）：检查
每个 $E_{m}\succeq0$ 且 $\sum_{m}E_{m}=\id$。(b) 计算在 $\ket{0}$ 上和在 $\ket{+}$ 上的三个结果
概率。(c) 解释为什么这不可能是投影测量，以及为什么结果为 $2$ 不能告诉你关于输入
量子比特的任何信息。
\end{exercise}

\begin{exercise}\label{ex:lec01-adder}
利用 \cref{sec:lec01-adder} 的半加器与全加器线路：(a) 对输入 $1+2$（即 $A_1A_0=01$ 且 $B_1B_0=10$）
追踪两比特加法器并确认
测得的输出是 $011_2=3$。(b) 数一数完整两比特加法器中 Toffoli 门和 $\CNOT$ 门的个数。
(c) 全加器会留下一个保存 $C_2$ 的草稿辅助比特，
同时输入保持完好；利用 Bennett 技巧解释，一旦
和位被复制出去，如何通过反计算中间的进位把任何草稿辅助比特重置回 $\ket{0}$，
以及当加法器作为更长计算中的子程序时为什么这很重要。
\end{exercise}

\readingsection
\begin{itemize}[leftmargin=1.4em]
  \item Lin Lin, \emph{Lecture Notes on Quantum Algorithms for Scientific
        Computation}, Ch.~1（公设、密度算符、量子线路、
        不可克隆、复杂度）--- \texttt{qasc\_notes.pdf}。
  \item Lin Lin, \emph{live notes}, Ch.~2（态、张量积、纠缠、
        混合态）--- \texttt{live\_notes\_0429.pdf}。
  \item A.~Childs, \emph{Lecture Notes on Quantum Algorithms}, Ch.~1--2（量子
        线路模型与通用性）--- \texttt{material/}。
  \item J.~Brody, \emph{The Joy of Quantum Computing}, Princeton University Press,
        2025, Ch.~4--7（隐形传态、不可克隆与不可超光速通信、CHSH/贝尔
        不等式、可逆算术）--- \texttt{material/}。
  \item 关于张量积以及计算基 $\ket{k}$ 记号（乘积态）的直观、循序渐进的
        讲解 ---
        \texttt{material/productstate\_zhihu.pdf}。
  \item M.~A.~Nielsen and I.~L.~Chuang, \emph{Quantum Computation and Quantum
        Information}, Ch.~1--2，提供完整的教科书式讲解。
\end{itemize}
